% paper56-analogy-transport.tex — Transporting Proofs Between Analogous Programs via Site Morphisms
%   Paper 56 of the Judgment Geometry series.
% Compile: pdflatex paper56-analogy-transport
\documentclass[11pt]{article}
\usepackage{lmodern}
% jugeo-common.tex — shared preamble for all Judgment Geometry papers
% Include via: % jugeo-common.tex — shared preamble for all Judgment Geometry papers
% Include via: % jugeo-common.tex — shared preamble for all Judgment Geometry papers
% Include via: \input{jugeo-common}

% ── Packages ────────────────────────────────────────────────────────────
\usepackage{amsmath,amssymb,amsthm,mathtools}
\usepackage{tikz-cd}
\usepackage{listings}
\usepackage{xcolor}
\usepackage{xspace}
\usepackage{booktabs}
\usepackage{multirow}
\usepackage{enumitem}
\usepackage[expansion=false]{microtype}
\usepackage{hyperref}
\usepackage{cleveref}
% \usepackage{thmtools}  % disabled: not available in this TeX installation
\usepackage{float}
\usepackage{caption}
\usepackage{subcaption}
\usepackage{tabularx}
\usepackage{stmaryrd}       % \llbracket, \rrbracket
\usepackage{bbm}            % \mathbbm{1}

% ── Page geometry ───────────────────────────────────────────────────────
\usepackage[margin=1in]{geometry}

% ── Theorem environments ────────────────────────────────────────────────
\theoremstyle{plain}
\newtheorem{theorem}{Theorem}[section]
\newtheorem{lemma}[theorem]{Lemma}
\newtheorem{proposition}[theorem]{Proposition}
\newtheorem{corollary}[theorem]{Corollary}

\theoremstyle{definition}
\newtheorem{definition}[theorem]{Definition}
\newtheorem{example}[theorem]{Example}
\newtheorem{construction}[theorem]{Construction}

\theoremstyle{remark}
\newtheorem{remark}[theorem]{Remark}
\newtheorem{notation}[theorem]{Notation}

% ── Python listings style ──────────────────────────────────────────────
\definecolor{jugeoBlue}{HTML}{2563EB}
\definecolor{jugeoGreen}{HTML}{059669}
\definecolor{jugeoRed}{HTML}{DC2626}
\definecolor{jugeoPurple}{HTML}{7C3AED}
\definecolor{jugeoGray}{HTML}{6B7280}
\definecolor{jugeoBg}{HTML}{F8FAFC}

\lstdefinestyle{jugeo-python}{
  language=Python,
  basicstyle=\ttfamily\small,
  keywordstyle=\color{jugeoBlue}\bfseries,
  stringstyle=\color{jugeoGreen},
  commentstyle=\color{jugeoGray}\itshape,
  numberstyle=\tiny\color{jugeoGray},
  backgroundcolor=\color{jugeoBg},
  numbers=left,
  numbersep=6pt,
  frame=single,
  framerule=0.4pt,
  rulecolor=\color{jugeoGray!40},
  breaklines=true,
  breakatwhitespace=true,
  showstringspaces=false,
  tabsize=4,
  xleftmargin=1.5em,
  framexleftmargin=1.5em,
  morekeywords={dataclass,frozen,field,Enum,IntEnum,auto,Any,
                Mapping,Sequence,Optional,match,case},
  literate={→}{{$\to$}}1 {←}{{$\leftarrow$}}1
           {≤}{{$\leq$}}1 {≥}{{$\geq$}}1 {≠}{{$\neq$}}1
           {∈}{{$\in$}}1 {∀}{{$\forall$}}1 {∃}{{$\exists$}}1
           {⊕}{{$\oplus$}}1 {⊖}{{$\ominus$}}1,
  escapeinside={(*@}{@*)},
}
\lstset{style=jugeo-python}

\lstdefinestyle{jugeo-output}{
  basicstyle=\ttfamily\small,
  backgroundcolor=\color{jugeoBg},
  frame=single,
  framerule=0.4pt,
  rulecolor=\color{jugeoGray!40},
  breaklines=true,
  showstringspaces=false,
  xleftmargin=1.5em,
  framexleftmargin=0.5em,
  numbers=none,
}

% ── JuGeo-specific macros ──────────────────────────────────────────────

% Judgment 8-tuple
\newcommand{\J}{\mathcal{J}}
\newcommand{\Jud}[8]{(#1,\,#2,\,#3,\,#4,\,#5,\,#6,\,#7,\,#8)}
\newcommand{\JudShort}[1]{\J_{#1}}

% Components
\newcommand{\coord}{\mathsf{c}}            % coordinate
\newcommand{\prop}{\varphi}                 % proposition
\newcommand{\carrier}{\mathcal{A}}          % carrier type
\newcommand{\evid}{\mathcal{E}}             % evidence bundle
\newcommand{\oblig}{\mathcal{O}}            % obligations
\newcommand{\obstr}{\mathcal{B}}            % obstructions
\newcommand{\trust}{\mathcal{T}}            % trust annotation
\newcommand{\prov}{\Pi}                     % provenance

% Trust algebra
\newcommand{\Talg}{\mathfrak{T}}            % trust ordered algebra
\newcommand{\Eadm}{\mathcal{E}_{\mathrm{adm}}}
\newcommand{\tleq}{\preceq}                 % trust ordering
\newcommand{\tjoin}{\oplus}                 % conservative join
\newcommand{\tatten}{\ominus}               % attenuation
\newcommand{\tpromote}{\uparrow_\pi}        % promotion
\newcommand{\tdemote}{\downarrow_\chi}       % demotion

% Trust tiers
\newcommand{\Tcontra}{\ensuremath{\mathsf{CONTRADICTED}}}
\newcommand{\Tunver}{\ensuremath{\mathsf{UNVERIFIED}}}
\newcommand{\Tcopilot}{\ensuremath{\mathsf{COPILOT}}}
\newcommand{\Toracle}{\ensuremath{\mathsf{ORACLE}}}
\newcommand{\Truntime}{\ensuremath{\mathsf{RUNTIME}}}
\newcommand{\Tsolver}{\ensuremath{\mathsf{SOLVER}}}
\newcommand{\Tproof}{\ensuremath{\mathsf{PROOF}}}

% Site and geometry
\newcommand{\Site}{\mathcal{S}}             % semantic site
\newcommand{\Cov}{\mathrm{Cov}}            % covering family
\newcommand{\Ob}{\mathrm{Ob}}              % objects
\newcommand{\Mor}{\mathrm{Mor}}            % morphisms
\newcommand{\res}{\mathrm{res}}            % restriction
\newcommand{\Cech}{\check{C}}              % Čech complex
\newcommand{\Hcoh}[1]{H^{#1}}             % cohomology group

% Descent
\newcommand{\descent}{\mathrm{desc}}
\newcommand{\glue}{\mathrm{glue}}
\newcommand{\overlap}[2]{#1 \cap #2}

% Semantic moves
\newcommand{\VERIFY}{\textsc{Verify}}
\newcommand{\CONSTRUCT}{\textsc{Construct}}
\newcommand{\NORMALIZE}{\textsc{Normalize}}
\newcommand{\GLUE}{\textsc{Glue}}
\newcommand{\RESTRICT}{\textsc{Restrict}}
\newcommand{\TRANSPORT}{\textsc{Transport}}
\newcommand{\REPAIR}{\textsc{Repair}}
\newcommand{\PROMOTE}{\textsc{Promote}}
\newcommand{\CHALLENGE}{\textsc{Challenge}}
\newcommand{\ESCALATE}{\textsc{Escalate}}

% Fragments
\newcommand{\QFLIA}{\texttt{QF\_LIA}}
\newcommand{\QFLRA}{\texttt{QF\_LRA}}
\newcommand{\QFBV}{\texttt{QF\_BV}}
\newcommand{\QFUF}{\texttt{QF\_UF}}

% Misc
\newcommand{\Python}{\textsc{Python}}
\newcommand{\JuGeo}{\textsc{JuGeo}}
\newcommand{\LEAN}{\textsc{Lean}\,4}
\newcommand{\Fstar}{F$^{\star}$}
\newcommand{\Coq}{\textsc{Coq}}
\newcommand{\Dafny}{\textsc{Dafny}}

% Common shortcuts
\newcommand{\ie}{i.e.\@\xspace}
\newcommand{\eg}{e.g.\@\xspace}
\newcommand{\cf}{cf.\@\xspace}
\newcommand{\wrt}{w.r.t.\@\xspace}

% Evaluation shortcuts
\newcommand{\numcases}{300}
\newcommand{\numspec}{100}
\newcommand{\numeq}{100}
\newcommand{\numbug}{100}
\newcommand{\medtime}{6.5\,ms}
\newcommand{\totaltime}{1.949\,s}


% ── Packages ────────────────────────────────────────────────────────────
\usepackage{amsmath,amssymb,amsthm,mathtools}
\usepackage{tikz-cd}
\usepackage{listings}
\usepackage{xcolor}
\usepackage{xspace}
\usepackage{booktabs}
\usepackage{multirow}
\usepackage{enumitem}
\usepackage[expansion=false]{microtype}
\usepackage{hyperref}
\usepackage{cleveref}
% \usepackage{thmtools}  % disabled: not available in this TeX installation
\usepackage{float}
\usepackage{caption}
\usepackage{subcaption}
\usepackage{tabularx}
\usepackage{stmaryrd}       % \llbracket, \rrbracket
\usepackage{bbm}            % \mathbbm{1}

% ── Page geometry ───────────────────────────────────────────────────────
\usepackage[margin=1in]{geometry}

% ── Theorem environments ────────────────────────────────────────────────
\theoremstyle{plain}
\newtheorem{theorem}{Theorem}[section]
\newtheorem{lemma}[theorem]{Lemma}
\newtheorem{proposition}[theorem]{Proposition}
\newtheorem{corollary}[theorem]{Corollary}

\theoremstyle{definition}
\newtheorem{definition}[theorem]{Definition}
\newtheorem{example}[theorem]{Example}
\newtheorem{construction}[theorem]{Construction}

\theoremstyle{remark}
\newtheorem{remark}[theorem]{Remark}
\newtheorem{notation}[theorem]{Notation}

% ── Python listings style ──────────────────────────────────────────────
\definecolor{jugeoBlue}{HTML}{2563EB}
\definecolor{jugeoGreen}{HTML}{059669}
\definecolor{jugeoRed}{HTML}{DC2626}
\definecolor{jugeoPurple}{HTML}{7C3AED}
\definecolor{jugeoGray}{HTML}{6B7280}
\definecolor{jugeoBg}{HTML}{F8FAFC}

\lstdefinestyle{jugeo-python}{
  language=Python,
  basicstyle=\ttfamily\small,
  keywordstyle=\color{jugeoBlue}\bfseries,
  stringstyle=\color{jugeoGreen},
  commentstyle=\color{jugeoGray}\itshape,
  numberstyle=\tiny\color{jugeoGray},
  backgroundcolor=\color{jugeoBg},
  numbers=left,
  numbersep=6pt,
  frame=single,
  framerule=0.4pt,
  rulecolor=\color{jugeoGray!40},
  breaklines=true,
  breakatwhitespace=true,
  showstringspaces=false,
  tabsize=4,
  xleftmargin=1.5em,
  framexleftmargin=1.5em,
  morekeywords={dataclass,frozen,field,Enum,IntEnum,auto,Any,
                Mapping,Sequence,Optional,match,case},
  literate={→}{{$\to$}}1 {←}{{$\leftarrow$}}1
           {≤}{{$\leq$}}1 {≥}{{$\geq$}}1 {≠}{{$\neq$}}1
           {∈}{{$\in$}}1 {∀}{{$\forall$}}1 {∃}{{$\exists$}}1
           {⊕}{{$\oplus$}}1 {⊖}{{$\ominus$}}1,
  escapeinside={(*@}{@*)},
}
\lstset{style=jugeo-python}

\lstdefinestyle{jugeo-output}{
  basicstyle=\ttfamily\small,
  backgroundcolor=\color{jugeoBg},
  frame=single,
  framerule=0.4pt,
  rulecolor=\color{jugeoGray!40},
  breaklines=true,
  showstringspaces=false,
  xleftmargin=1.5em,
  framexleftmargin=0.5em,
  numbers=none,
}

% ── JuGeo-specific macros ──────────────────────────────────────────────

% Judgment 8-tuple
\newcommand{\J}{\mathcal{J}}
\newcommand{\Jud}[8]{(#1,\,#2,\,#3,\,#4,\,#5,\,#6,\,#7,\,#8)}
\newcommand{\JudShort}[1]{\J_{#1}}

% Components
\newcommand{\coord}{\mathsf{c}}            % coordinate
\newcommand{\prop}{\varphi}                 % proposition
\newcommand{\carrier}{\mathcal{A}}          % carrier type
\newcommand{\evid}{\mathcal{E}}             % evidence bundle
\newcommand{\oblig}{\mathcal{O}}            % obligations
\newcommand{\obstr}{\mathcal{B}}            % obstructions
\newcommand{\trust}{\mathcal{T}}            % trust annotation
\newcommand{\prov}{\Pi}                     % provenance

% Trust algebra
\newcommand{\Talg}{\mathfrak{T}}            % trust ordered algebra
\newcommand{\Eadm}{\mathcal{E}_{\mathrm{adm}}}
\newcommand{\tleq}{\preceq}                 % trust ordering
\newcommand{\tjoin}{\oplus}                 % conservative join
\newcommand{\tatten}{\ominus}               % attenuation
\newcommand{\tpromote}{\uparrow_\pi}        % promotion
\newcommand{\tdemote}{\downarrow_\chi}       % demotion

% Trust tiers
\newcommand{\Tcontra}{\ensuremath{\mathsf{CONTRADICTED}}}
\newcommand{\Tunver}{\ensuremath{\mathsf{UNVERIFIED}}}
\newcommand{\Tcopilot}{\ensuremath{\mathsf{COPILOT}}}
\newcommand{\Toracle}{\ensuremath{\mathsf{ORACLE}}}
\newcommand{\Truntime}{\ensuremath{\mathsf{RUNTIME}}}
\newcommand{\Tsolver}{\ensuremath{\mathsf{SOLVER}}}
\newcommand{\Tproof}{\ensuremath{\mathsf{PROOF}}}

% Site and geometry
\newcommand{\Site}{\mathcal{S}}             % semantic site
\newcommand{\Cov}{\mathrm{Cov}}            % covering family
\newcommand{\Ob}{\mathrm{Ob}}              % objects
\newcommand{\Mor}{\mathrm{Mor}}            % morphisms
\newcommand{\res}{\mathrm{res}}            % restriction
\newcommand{\Cech}{\check{C}}              % Čech complex
\newcommand{\Hcoh}[1]{H^{#1}}             % cohomology group

% Descent
\newcommand{\descent}{\mathrm{desc}}
\newcommand{\glue}{\mathrm{glue}}
\newcommand{\overlap}[2]{#1 \cap #2}

% Semantic moves
\newcommand{\VERIFY}{\textsc{Verify}}
\newcommand{\CONSTRUCT}{\textsc{Construct}}
\newcommand{\NORMALIZE}{\textsc{Normalize}}
\newcommand{\GLUE}{\textsc{Glue}}
\newcommand{\RESTRICT}{\textsc{Restrict}}
\newcommand{\TRANSPORT}{\textsc{Transport}}
\newcommand{\REPAIR}{\textsc{Repair}}
\newcommand{\PROMOTE}{\textsc{Promote}}
\newcommand{\CHALLENGE}{\textsc{Challenge}}
\newcommand{\ESCALATE}{\textsc{Escalate}}

% Fragments
\newcommand{\QFLIA}{\texttt{QF\_LIA}}
\newcommand{\QFLRA}{\texttt{QF\_LRA}}
\newcommand{\QFBV}{\texttt{QF\_BV}}
\newcommand{\QFUF}{\texttt{QF\_UF}}

% Misc
\newcommand{\Python}{\textsc{Python}}
\newcommand{\JuGeo}{\textsc{JuGeo}}
\newcommand{\LEAN}{\textsc{Lean}\,4}
\newcommand{\Fstar}{F$^{\star}$}
\newcommand{\Coq}{\textsc{Coq}}
\newcommand{\Dafny}{\textsc{Dafny}}

% Common shortcuts
\newcommand{\ie}{i.e.\@\xspace}
\newcommand{\eg}{e.g.\@\xspace}
\newcommand{\cf}{cf.\@\xspace}
\newcommand{\wrt}{w.r.t.\@\xspace}

% Evaluation shortcuts
\newcommand{\numcases}{300}
\newcommand{\numspec}{100}
\newcommand{\numeq}{100}
\newcommand{\numbug}{100}
\newcommand{\medtime}{6.5\,ms}
\newcommand{\totaltime}{1.949\,s}


% ── Packages ────────────────────────────────────────────────────────────
\usepackage{amsmath,amssymb,amsthm,mathtools}
\usepackage{tikz-cd}
\usepackage{listings}
\usepackage{xcolor}
\usepackage{xspace}
\usepackage{booktabs}
\usepackage{multirow}
\usepackage{enumitem}
\usepackage[expansion=false]{microtype}
\usepackage{hyperref}
\usepackage{cleveref}
% \usepackage{thmtools}  % disabled: not available in this TeX installation
\usepackage{float}
\usepackage{caption}
\usepackage{subcaption}
\usepackage{tabularx}
\usepackage{stmaryrd}       % \llbracket, \rrbracket
\usepackage{bbm}            % \mathbbm{1}

% ── Page geometry ───────────────────────────────────────────────────────
\usepackage[margin=1in]{geometry}

% ── Theorem environments ────────────────────────────────────────────────
\theoremstyle{plain}
\newtheorem{theorem}{Theorem}[section]
\newtheorem{lemma}[theorem]{Lemma}
\newtheorem{proposition}[theorem]{Proposition}
\newtheorem{corollary}[theorem]{Corollary}

\theoremstyle{definition}
\newtheorem{definition}[theorem]{Definition}
\newtheorem{example}[theorem]{Example}
\newtheorem{construction}[theorem]{Construction}

\theoremstyle{remark}
\newtheorem{remark}[theorem]{Remark}
\newtheorem{notation}[theorem]{Notation}

% ── Python listings style ──────────────────────────────────────────────
\definecolor{jugeoBlue}{HTML}{2563EB}
\definecolor{jugeoGreen}{HTML}{059669}
\definecolor{jugeoRed}{HTML}{DC2626}
\definecolor{jugeoPurple}{HTML}{7C3AED}
\definecolor{jugeoGray}{HTML}{6B7280}
\definecolor{jugeoBg}{HTML}{F8FAFC}

\lstdefinestyle{jugeo-python}{
  language=Python,
  basicstyle=\ttfamily\small,
  keywordstyle=\color{jugeoBlue}\bfseries,
  stringstyle=\color{jugeoGreen},
  commentstyle=\color{jugeoGray}\itshape,
  numberstyle=\tiny\color{jugeoGray},
  backgroundcolor=\color{jugeoBg},
  numbers=left,
  numbersep=6pt,
  frame=single,
  framerule=0.4pt,
  rulecolor=\color{jugeoGray!40},
  breaklines=true,
  breakatwhitespace=true,
  showstringspaces=false,
  tabsize=4,
  xleftmargin=1.5em,
  framexleftmargin=1.5em,
  morekeywords={dataclass,frozen,field,Enum,IntEnum,auto,Any,
                Mapping,Sequence,Optional,match,case},
  literate={→}{{$\to$}}1 {←}{{$\leftarrow$}}1
           {≤}{{$\leq$}}1 {≥}{{$\geq$}}1 {≠}{{$\neq$}}1
           {∈}{{$\in$}}1 {∀}{{$\forall$}}1 {∃}{{$\exists$}}1
           {⊕}{{$\oplus$}}1 {⊖}{{$\ominus$}}1,
  escapeinside={(*@}{@*)},
}
\lstset{style=jugeo-python}

\lstdefinestyle{jugeo-output}{
  basicstyle=\ttfamily\small,
  backgroundcolor=\color{jugeoBg},
  frame=single,
  framerule=0.4pt,
  rulecolor=\color{jugeoGray!40},
  breaklines=true,
  showstringspaces=false,
  xleftmargin=1.5em,
  framexleftmargin=0.5em,
  numbers=none,
}

% ── JuGeo-specific macros ──────────────────────────────────────────────

% Judgment 8-tuple
\newcommand{\J}{\mathcal{J}}
\newcommand{\Jud}[8]{(#1,\,#2,\,#3,\,#4,\,#5,\,#6,\,#7,\,#8)}
\newcommand{\JudShort}[1]{\J_{#1}}

% Components
\newcommand{\coord}{\mathsf{c}}            % coordinate
\newcommand{\prop}{\varphi}                 % proposition
\newcommand{\carrier}{\mathcal{A}}          % carrier type
\newcommand{\evid}{\mathcal{E}}             % evidence bundle
\newcommand{\oblig}{\mathcal{O}}            % obligations
\newcommand{\obstr}{\mathcal{B}}            % obstructions
\newcommand{\trust}{\mathcal{T}}            % trust annotation
\newcommand{\prov}{\Pi}                     % provenance

% Trust algebra
\newcommand{\Talg}{\mathfrak{T}}            % trust ordered algebra
\newcommand{\Eadm}{\mathcal{E}_{\mathrm{adm}}}
\newcommand{\tleq}{\preceq}                 % trust ordering
\newcommand{\tjoin}{\oplus}                 % conservative join
\newcommand{\tatten}{\ominus}               % attenuation
\newcommand{\tpromote}{\uparrow_\pi}        % promotion
\newcommand{\tdemote}{\downarrow_\chi}       % demotion

% Trust tiers
\newcommand{\Tcontra}{\ensuremath{\mathsf{CONTRADICTED}}}
\newcommand{\Tunver}{\ensuremath{\mathsf{UNVERIFIED}}}
\newcommand{\Tcopilot}{\ensuremath{\mathsf{COPILOT}}}
\newcommand{\Toracle}{\ensuremath{\mathsf{ORACLE}}}
\newcommand{\Truntime}{\ensuremath{\mathsf{RUNTIME}}}
\newcommand{\Tsolver}{\ensuremath{\mathsf{SOLVER}}}
\newcommand{\Tproof}{\ensuremath{\mathsf{PROOF}}}

% Site and geometry
\newcommand{\Site}{\mathcal{S}}             % semantic site
\newcommand{\Cov}{\mathrm{Cov}}            % covering family
\newcommand{\Ob}{\mathrm{Ob}}              % objects
\newcommand{\Mor}{\mathrm{Mor}}            % morphisms
\newcommand{\res}{\mathrm{res}}            % restriction
\newcommand{\Cech}{\check{C}}              % Čech complex
\newcommand{\Hcoh}[1]{H^{#1}}             % cohomology group

% Descent
\newcommand{\descent}{\mathrm{desc}}
\newcommand{\glue}{\mathrm{glue}}
\newcommand{\overlap}[2]{#1 \cap #2}

% Semantic moves
\newcommand{\VERIFY}{\textsc{Verify}}
\newcommand{\CONSTRUCT}{\textsc{Construct}}
\newcommand{\NORMALIZE}{\textsc{Normalize}}
\newcommand{\GLUE}{\textsc{Glue}}
\newcommand{\RESTRICT}{\textsc{Restrict}}
\newcommand{\TRANSPORT}{\textsc{Transport}}
\newcommand{\REPAIR}{\textsc{Repair}}
\newcommand{\PROMOTE}{\textsc{Promote}}
\newcommand{\CHALLENGE}{\textsc{Challenge}}
\newcommand{\ESCALATE}{\textsc{Escalate}}

% Fragments
\newcommand{\QFLIA}{\texttt{QF\_LIA}}
\newcommand{\QFLRA}{\texttt{QF\_LRA}}
\newcommand{\QFBV}{\texttt{QF\_BV}}
\newcommand{\QFUF}{\texttt{QF\_UF}}

% Misc
\newcommand{\Python}{\textsc{Python}}
\newcommand{\JuGeo}{\textsc{JuGeo}}
\newcommand{\LEAN}{\textsc{Lean}\,4}
\newcommand{\Fstar}{F$^{\star}$}
\newcommand{\Coq}{\textsc{Coq}}
\newcommand{\Dafny}{\textsc{Dafny}}

% Common shortcuts
\newcommand{\ie}{i.e.\@\xspace}
\newcommand{\eg}{e.g.\@\xspace}
\newcommand{\cf}{cf.\@\xspace}
\newcommand{\wrt}{w.r.t.\@\xspace}

% Evaluation shortcuts
\newcommand{\numcases}{300}
\newcommand{\numspec}{100}
\newcommand{\numeq}{100}
\newcommand{\numbug}{100}
\newcommand{\medtime}{6.5\,ms}
\newcommand{\totaltime}{1.949\,s}

% experiment-data.tex — AUTO-GENERATED from real jugeo CLI runs
% DO NOT EDIT — regenerate with: python3 experiments/run_universal_metrics.py
% Generated from 30 programs across 6 domains

\newcommand{\expTotalPrograms}{30}
\newcommand{\expVerified}{30}
\newcommand{\expAccuracy}{100.0\%}
\newcommand{\expCoordsMin}{2}
\newcommand{\expCoordsMax}{10}
\newcommand{\expCoordsMean}{5.2}
\newcommand{\expCoordsMedian}{5.5}
\newcommand{\expPropsMin}{4}
\newcommand{\expPropsMax}{20}
\newcommand{\expPropsMean}{10.4}
\newcommand{\expPropsSum}{311}
\newcommand{\expPropsOkSum}{310}
\newcommand{\expTimeMin}{3.506\,s}
\newcommand{\expTimeMax}{6.714\,s}
\newcommand{\expTimeMean}{4.64\,s}
\newcommand{\expTimeMedian}{4.508\,s}
\newcommand{\expTimeTotal}{139.204\,s}
\newcommand{\expObstructionTotal}{0}
\newcommand{\expHOneZero}{30}
\newcommand{\expDomainCount}{6}
\newcommand{\expTrustCOPILOTSUGGESTED}{30}
\newcommand{\expDomainDsCount}{5}
\newcommand{\expDomainDsVerified}{5}
\newcommand{\expDomainDsAvgCoords}{7.6}
\newcommand{\expDomainDsAvgProps}{15.2}
\newcommand{\expDomainMathCount}{5}
\newcommand{\expDomainMathVerified}{5}
\newcommand{\expDomainMathAvgCoords}{4.8}
\newcommand{\expDomainMathAvgProps}{9.4}
\newcommand{\expDomainSmCount}{5}
\newcommand{\expDomainSmVerified}{5}
\newcommand{\expDomainSmAvgCoords}{7.6}
\newcommand{\expDomainSmAvgProps}{15.2}
\newcommand{\expDomainSortCount}{5}
\newcommand{\expDomainSortVerified}{5}
\newcommand{\expDomainSortAvgCoords}{2.4}
\newcommand{\expDomainSortAvgProps}{4.8}
\newcommand{\expDomainStrCount}{5}
\newcommand{\expDomainStrVerified}{5}
\newcommand{\expDomainStrAvgCoords}{3.2}
\newcommand{\expDomainStrAvgProps}{6.4}
\newcommand{\expDomainWebCount}{5}
\newcommand{\expDomainWebVerified}{5}
\newcommand{\expDomainWebAvgCoords}{5.6}
\newcommand{\expDomainWebAvgProps}{11.2}

% subsystem-data.tex — AUTO-GENERATED from real jugeo subsystem experiments
% DO NOT EDIT — regenerate with: python3 experiments/run_subsystem_experiments.py

\newcommand{\subCliTotal}{10}
\newcommand{\subCliVerified}{9}
\newcommand{\subCliAccuracy}{90.0\%}
\newcommand{\subCliMeanTime}{4.132\,s}
\newcommand{\subCliMinTime}{3.472\,s}
\newcommand{\subCliMaxTime}{5.372\,s}
\newcommand{\subCliTotalCoords}{54}
\newcommand{\subCliTotalProps}{107}
\newcommand{\subCliTotalPropsOk}{106}
\newcommand{\subSiteCalls}{90}
\newcommand{\subSiteSuccessful}{69}
\newcommand{\subSiteSuccessRate}{76.7\%}
\newcommand{\subSiteMeanTime}{0.0001\,s}
\newcommand{\subSiteTotalTime}{0.005\,s}
\newcommand{\subSpecificationsatisfactionSmallTime}{0.0003\,s}
\newcommand{\subSpecificationsatisfactionMediumTime}{0.0001\,s}
\newcommand{\subSpecificationsatisfactionLargeTime}{0.0001\,s}
\newcommand{\subInhabitantfleetSmallTime}{0.0\,s}
\newcommand{\subInhabitantfleetMediumTime}{0.0\,s}
\newcommand{\subInhabitantfleetLargeTime}{0.0\,s}
\newcommand{\subTheoremecologySmallTime}{0.0\,s}
\newcommand{\subTheoremecologyMediumTime}{0.0\,s}
\newcommand{\subTheoremecologyLargeTime}{0.0\,s}
\newcommand{\subStatespaceexplorationSmallTime}{0.0\,s}
\newcommand{\subStatespaceexplorationMediumTime}{0.0\,s}
\newcommand{\subStatespaceexplorationLargeTime}{0.0\,s}
\newcommand{\subRepairsemanticsSmallTime}{0.0\,s}
\newcommand{\subRepairsemanticsMediumTime}{0.0\,s}
\newcommand{\subRepairsemanticsLargeTime}{0.0\,s}
\newcommand{\subReplaygluingSmallTime}{0.0\,s}
\newcommand{\subReplaygluingMediumTime}{0.0\,s}
\newcommand{\subReplaygluingLargeTime}{0.0\,s}
\newcommand{\subSemanticfuturesSmallTime}{0.0\,s}
\newcommand{\subSemanticfuturesMediumTime}{0.0\,s}
\newcommand{\subSemanticfuturesLargeTime}{0.0\,s}
\newcommand{\subHypercovertreatySmallTime}{0.0\,s}
\newcommand{\subHypercovertreatyMediumTime}{0.0\,s}
\newcommand{\subHypercovertreatyLargeTime}{0.0\,s}
\newcommand{\subEncodeforsolverSmallTime}{0.0026\,s}
\newcommand{\subEncodeforsolverMediumTime}{0.0002\,s}
\newcommand{\subEncodeforsolverLargeTime}{0.0001\,s}
\newcommand{\subInterfaceroutingSmallTime}{0.0\,s}
\newcommand{\subInterfaceroutingMediumTime}{0.0\,s}
\newcommand{\subInterfaceroutingLargeTime}{0.0\,s}
\newcommand{\subOrchestrateverificationSmallTime}{0.0002\,s}
\newcommand{\subOrchestrateverificationMediumTime}{0.0\,s}
\newcommand{\subOrchestrateverificationLargeTime}{0.0\,s}
\newcommand{\subRunfulldescentSmallTime}{0.0\,s}
\newcommand{\subRunfulldescentMediumTime}{0.0\,s}
\newcommand{\subRunfulldescentLargeTime}{0.0\,s}
\newcommand{\subKernellifecycleSmallTime}{0.0\,s}
\newcommand{\subKernellifecycleMediumTime}{0.0\,s}
\newcommand{\subKernellifecycleLargeTime}{0.0\,s}
\newcommand{\subDiscoverypipelineSmallTime}{0.0\,s}
\newcommand{\subDiscoverypipelineMediumTime}{0.0\,s}
\newcommand{\subDiscoverypipelineLargeTime}{0.0\,s}
\newcommand{\subAnalogytransportSmallTime}{0.0\,s}
\newcommand{\subAnalogytransportMediumTime}{0.0\,s}
\newcommand{\subAnalogytransportLargeTime}{0.0\,s}
\newcommand{\subFormalcoresiteSmallTime}{0.0001\,s}
\newcommand{\subFormalcoresiteMediumTime}{0.0\,s}
\newcommand{\subFormalcoresiteLargeTime}{0.0\,s}
\newcommand{\subProblematlasSmallTime}{0.0\,s}
\newcommand{\subProblematlasMediumTime}{0.0\,s}
\newcommand{\subProblematlasLargeTime}{0.0\,s}
\newcommand{\subPublicalignmentSmallTime}{0.0\,s}
\newcommand{\subPublicalignmentMediumTime}{0.0\,s}
\newcommand{\subPublicalignmentLargeTime}{0.0\,s}
\newcommand{\subRegimebootstrappingSmallTime}{0.0\,s}
\newcommand{\subRegimebootstrappingMediumTime}{0.0\,s}
\newcommand{\subRegimebootstrappingLargeTime}{0.0\,s}
\newcommand{\subRelationalrefinementSmallTime}{0.0\,s}
\newcommand{\subRelationalrefinementMediumTime}{0.0\,s}
\newcommand{\subRelationalrefinementLargeTime}{0.0\,s}
\newcommand{\subSemanticclosureSmallTime}{0.0\,s}
\newcommand{\subSemanticclosureMediumTime}{0.0\,s}
\newcommand{\subSemanticclosureLargeTime}{0.0\,s}
\newcommand{\subTheoremeconomicsSmallTime}{0.0001\,s}
\newcommand{\subTheoremeconomicsMediumTime}{0.0\,s}
\newcommand{\subTheoremeconomicsLargeTime}{0.0\,s}
\newcommand{\subBenchmarksuiteSmallTime}{0.0\,s}
\newcommand{\subBenchmarksuiteMediumTime}{0.0\,s}
\newcommand{\subBenchmarksuiteLargeTime}{0.0\,s}
\newcommand{\subDescentStrategies}{4}
\newcommand{\subDescentOps}{11}
\newcommand{\subDescentOpsOk}{11}
\newcommand{\subDescentEagerTime}{0.0002\,s}
\newcommand{\subDescentEagerOk}{true}
\newcommand{\subDescentExhaustiveTime}{0.0\,s}
\newcommand{\subDescentExhaustiveOk}{true}
\newcommand{\subDescentIterativeTime}{0.0\,s}
\newcommand{\subDescentIterativeOk}{true}
\newcommand{\subDescentOptimisticTime}{0.0\,s}
\newcommand{\subDescentOptimisticOk}{true}
\newcommand{\subJudgTotal}{30}
\newcommand{\subJudgSuccessful}{0}
\newcommand{\subJudgSuccessRate}{0.0\%}
\newcommand{\subJudgMeanTimeUs}{7.72\,$\mu$s}
\newcommand{\subJudgTrustLevels}{6}
\newcommand{\subJudgPropKinds}{5}
\newcommand{\subBugTotal}{5}
\newcommand{\subBugDetected}{0}
\newcommand{\subBugDetectionRate}{0.0\%}
\newcommand{\subBugFalsePositiveRate}{0.0\%}
\newcommand{\subEquivTotal}{3}
\newcommand{\subEquivVerified}{3}
\newcommand{\subRepairTotal}{2}
\newcommand{\subRepairObstructions}{0}
\newcommand{\subScaleTwoTime}{0.0001\,s}
\newcommand{\subScaleFourTime}{0.0\,s}
\newcommand{\subScaleEightTime}{0.0\,s}
\newcommand{\subScaleTwelveTime}{0.0\,s}
\newcommand{\subScaleSixteenTime}{0.0\,s}
\newcommand{\subScaleTwentyTime}{0.0\,s}
\newcommand{\subTotalExperimentTime}{95.6\,s}

% comprehensive-data.tex — AUTO-GENERATED from real jugeo experiments
% DO NOT EDIT — regenerate with: python3 experiments/run_comprehensive_experiments.py
% Generated: 2026-03-23 09:19:06

% --- Size scaling metrics ---
\newcommand{\compTotalPrograms}{18}
\newcommand{\compTotalVerified}{18}
\newcommand{\compOverallAccuracy}{100.0\%}
\newcommand{\compTinyCount}{5}
\newcommand{\compTinyVerified}{5}
\newcommand{\compTinyMeanCoords}{3}
\newcommand{\compTinyMeanProps}{6}
\newcommand{\compTinyMeanTime}{4.95\,s}
\newcommand{\compTinyMeanLines}{5}
\newcommand{\compTinyMinTime}{4.45\,s}
\newcommand{\compTinyMaxTime}{5.51\,s}
\newcommand{\compSmallCount}{5}
\newcommand{\compSmallVerified}{5}
\newcommand{\compSmallMeanCoords}{3.6}
\newcommand{\compSmallMeanProps}{7.2}
\newcommand{\compSmallMeanTime}{4.85\,s}
\newcommand{\compSmallMeanLines}{16.0}
\newcommand{\compSmallMinTime}{4.30\,s}
\newcommand{\compSmallMaxTime}{5.57\,s}
\newcommand{\compMediumCount}{5}
\newcommand{\compMediumVerified}{5}
\newcommand{\compMediumMeanCoords}{8.6}
\newcommand{\compMediumMeanProps}{17.2}
\newcommand{\compMediumMeanTime}{4.47\,s}
\newcommand{\compMediumMeanLines}{45.0}
\newcommand{\compMediumMinTime}{4.26\,s}
\newcommand{\compMediumMaxTime}{4.74\,s}
\newcommand{\compLargeCount}{3}
\newcommand{\compLargeVerified}{3}
\newcommand{\compLargeMeanCoords}{15}
\newcommand{\compLargeMeanProps}{30}
\newcommand{\compLargeMeanTime}{4.31\,s}
\newcommand{\compLargeMeanLines}{79.0}
\newcommand{\compLargeMinTime}{4.27\,s}
\newcommand{\compLargeMaxTime}{4.38\,s}
% --- Aggregate timing ---
\newcommand{\compTimeMean}{4.68\,s}
\newcommand{\compTimeMedian}{4.50\,s}
\newcommand{\compTimeMin}{4.265\,s}
\newcommand{\compTimeMax}{5.571\,s}
\newcommand{\compTimeTotal}{84.3\,s}
\newcommand{\compTimeStdev}{0.45\,s}
% --- Aggregate structure ---
\newcommand{\compCoordsMean}{6.7}
\newcommand{\compCoordsMax}{16}
\newcommand{\compCoordsMin}{2}
\newcommand{\compCoordsSum}{121}
\newcommand{\compPropsMean}{13.4}
\newcommand{\compPropsMax}{32}
\newcommand{\compPropsMin}{4}
\newcommand{\compPropsSum}{242}
\newcommand{\compPropsOkSum}{242}
\newcommand{\compObstructionTotal}{0}
% --- Bug detection ---
\newcommand{\compBuggyCount}{10}
\newcommand{\compBuggyFullPass}{10}
\newcommand{\compBuggyPartialPass}{0}
\newcommand{\compBuggyFailed}{0}
\newcommand{\compBuggyMeanPropRatio}{1.00}
\newcommand{\compCorrectMeanPropRatio}{1.00}
\newcommand{\compBuggyMeanTime}{5.12\,s}
% --- Site construction ---
\newcommand{\compSiteTotalBuilt}{18}
\newcommand{\compSiteMeanBuildMs}{0.05}
\newcommand{\compSiteMaxBuildMs}{0.14}
\newcommand{\compSiteMeanCoords}{6.5}
\newcommand{\compSiteMeanMorphisms}{14.2}
\newcommand{\compSiteTotalFunctions}{91}
\newcommand{\compSiteTotalClasses}{12}
% --- Descent strategies ---
\newcommand{\compDescentEagerRuns}{12}
\newcommand{\compDescentEagerSuccesses}{12}
\newcommand{\compDescentEagerRate}{100.0\%}
\newcommand{\compDescentEagerMeanMs}{0.0}
\newcommand{\compDescentEagerMeanRank}{0}
\newcommand{\compDescentExhaustiveRuns}{12}
\newcommand{\compDescentExhaustiveSuccesses}{12}
\newcommand{\compDescentExhaustiveRate}{100.0\%}
\newcommand{\compDescentExhaustiveMeanMs}{0.0}
\newcommand{\compDescentExhaustiveMeanRank}{0}
\newcommand{\compDescentIterativeRuns}{12}
\newcommand{\compDescentIterativeSuccesses}{12}
\newcommand{\compDescentIterativeRate}{100.0\%}
\newcommand{\compDescentIterativeMeanMs}{0.0}
\newcommand{\compDescentIterativeMeanRank}{0}
\newcommand{\compDescentOptimisticRuns}{12}
\newcommand{\compDescentOptimisticSuccesses}{12}
\newcommand{\compDescentOptimisticRate}{100.0\%}
\newcommand{\compDescentOptimisticMeanMs}{0.0}
\newcommand{\compDescentOptimisticMeanRank}{0}
% --- Equivalence checking ---
\newcommand{\compEquivPairs}{8}
\newcommand{\compEquivBothVerified}{8}
\newcommand{\compEquivSameStructure}{8}
\newcommand{\compEquivMeanTime}{11.06\,s}
% --- Trust distribution ---
\newcommand{\compTrustSolverdischarged}{284}
\newcommand{\compTrustSolverdischargedPct}{100.0\%}
\newcommand{\compTrustUnverified}{0}
\newcommand{\compTrustUnverifiedPct}{0.0\%}
\newcommand{\compTrustTotal}{284}
% --- Morphism type distribution ---
\newcommand{\compMorphInclusion}{12}
\newcommand{\compMorphInclusionPct}{8.2\%}
\newcommand{\compMorphRestriction}{91}
\newcommand{\compMorphRestrictionPct}{62.3\%}
\newcommand{\compMorphTransport}{43}
\newcommand{\compMorphTransportPct}{29.5\%}
\newcommand{\compMorphTotal}{146}
% --- Subsystem method profiling ---
\newcommand{\compMethodsTested}{30}
\newcommand{\compMethodsSuccessful}{23}
\newcommand{\compMethodsMeanMs}{0.02}
\newcommand{\compProfAnalogyTransportMeanMs}{0.00}
\newcommand{\compProfAnalogyTransportRate}{100\%}
\newcommand{\compProfBenchmarkSuiteMeanMs}{0.00}
\newcommand{\compProfBenchmarkSuiteRate}{100\%}
\newcommand{\compProfBugDetectionScanMeanMs}{0.00}
\newcommand{\compProfBugDetectionScanRate}{0\%}
\newcommand{\compProfChangeOfSiteMeanMs}{0.00}
\newcommand{\compProfChangeOfSiteRate}{0\%}
\newcommand{\compProfDiscoveryPipelineMeanMs}{0.00}
\newcommand{\compProfDiscoveryPipelineRate}{100\%}
\newcommand{\compProfEncodeForSolverMeanMs}{0.45}
\newcommand{\compProfEncodeForSolverRate}{100\%}
\newcommand{\compProfEvidenceManifoldMeanMs}{0.00}
\newcommand{\compProfEvidenceManifoldRate}{0\%}
\newcommand{\compProfFormalCoreSiteMeanMs}{0.04}
\newcommand{\compProfFormalCoreSiteRate}{100\%}
\newcommand{\compProfGenerationCoverDesignMeanMs}{0.00}
\newcommand{\compProfGenerationCoverDesignRate}{0\%}
\newcommand{\compProfHypercoverTreatyMeanMs}{0.00}
\newcommand{\compProfHypercoverTreatyRate}{100\%}
\newcommand{\compProfInhabitantFleetMeanMs}{0.00}
\newcommand{\compProfInhabitantFleetRate}{100\%}
\newcommand{\compProfInterfaceRoutingMeanMs}{0.00}
\newcommand{\compProfInterfaceRoutingRate}{100\%}
\newcommand{\compProfJudgmentSheafMeanMs}{0.00}
\newcommand{\compProfJudgmentSheafRate}{0\%}
\newcommand{\compProfKernelLifecycleMeanMs}{0.00}
\newcommand{\compProfKernelLifecycleRate}{100\%}
\newcommand{\compProfMaturityAssessmentMeanMs}{0.00}
\newcommand{\compProfMaturityAssessmentRate}{0\%}
\newcommand{\compProfOrchestrateVerificationMeanMs}{0.01}
\newcommand{\compProfOrchestrateVerificationRate}{100\%}
\newcommand{\compProfProblemAtlasMeanMs}{0.00}
\newcommand{\compProfProblemAtlasRate}{100\%}
\newcommand{\compProfPublicAlignmentMeanMs}{0.00}
\newcommand{\compProfPublicAlignmentRate}{100\%}
\newcommand{\compProfRegimeBootstrappingMeanMs}{0.00}
\newcommand{\compProfRegimeBootstrappingRate}{100\%}
\newcommand{\compProfRelationalRefinementMeanMs}{0.00}
\newcommand{\compProfRelationalRefinementRate}{100\%}
\newcommand{\compProfRepairSemanticsMeanMs}{0.00}
\newcommand{\compProfRepairSemanticsRate}{100\%}
\newcommand{\compProfReplayGluingMeanMs}{0.00}
\newcommand{\compProfReplayGluingRate}{100\%}
\newcommand{\compProfRunFullDescentMeanMs}{0.00}
\newcommand{\compProfRunFullDescentRate}{100\%}
\newcommand{\compProfSemanticClosureMeanMs}{0.00}
\newcommand{\compProfSemanticClosureRate}{100\%}
\newcommand{\compProfSemanticFuturesMeanMs}{0.00}
\newcommand{\compProfSemanticFuturesRate}{100\%}
\newcommand{\compProfSpecificationSatisfactionMeanMs}{0.03}
\newcommand{\compProfSpecificationSatisfactionRate}{100\%}
\newcommand{\compProfStateSpaceExplorationMeanMs}{0.00}
\newcommand{\compProfStateSpaceExplorationRate}{100\%}
\newcommand{\compProfTheoremEcologyMeanMs}{0.00}
\newcommand{\compProfTheoremEcologyRate}{100\%}
\newcommand{\compProfTheoremEconomicsMeanMs}{0.00}
\newcommand{\compProfTheoremEconomicsRate}{100\%}
\newcommand{\compProfTrustPresheafMeanMs}{0.00}
\newcommand{\compProfTrustPresheafRate}{0\%}

% supplementary-data.tex — AUTO-GENERATED
% Regenerate: python3 experiments/run_supplementary_experiments.py
% Generated: 2026-03-23 09:57:40

\newcommand{\suppDomainSortAvgTime}{3.59\,s}
\newcommand{\suppDomainSortMinTime}{3.18\,s}
\newcommand{\suppDomainSortMaxTime}{4.00\,s}
\newcommand{\suppDomainDsAvgTime}{3.41\,s}
\newcommand{\suppDomainDsMinTime}{3.27\,s}
\newcommand{\suppDomainDsMaxTime}{3.75\,s}
\newcommand{\suppDomainMathAvgTime}{3.76\,s}
\newcommand{\suppDomainMathMinTime}{3.15\,s}
\newcommand{\suppDomainMathMaxTime}{4.05\,s}
\newcommand{\suppDomainStrAvgTime}{3.27\,s}
\newcommand{\suppDomainStrMinTime}{3.05\,s}
\newcommand{\suppDomainStrMaxTime}{3.47\,s}
\newcommand{\suppDomainWebAvgTime}{3.33\,s}
\newcommand{\suppDomainWebMinTime}{3.25\,s}
\newcommand{\suppDomainWebMaxTime}{3.47\,s}
\newcommand{\suppDomainSmAvgTime}{3.81\,s}
\newcommand{\suppDomainSmMinTime}{3.41\,s}
\newcommand{\suppDomainSmMaxTime}{4.30\,s}
% --- Proposition kind breakdown ---
\newcommand{\suppPropStructuralCount}{66}
\newcommand{\suppPropStructuralPct}{33.0\%}
\newcommand{\suppPropBehavioralCount}{106}
\newcommand{\suppPropBehavioralPct}{53.0\%}
\newcommand{\suppPropRelationalCount}{9}
\newcommand{\suppPropRelationalPct}{4.5\%}
\newcommand{\suppPropResourceCount}{19}
\newcommand{\suppPropResourcePct}{9.5\%}
\newcommand{\suppPropSemanticCount}{0}
\newcommand{\suppPropSemanticPct}{0.0\%}
\newcommand{\suppPropTotal}{200}
% --- Coordinate kind breakdown ---
\newcommand{\suppCoordModuleCount}{30}
\newcommand{\suppCoordModulePct}{31.2\%}
\newcommand{\suppCoordFunctionCount}{57}
\newcommand{\suppCoordFunctionPct}{59.4\%}
\newcommand{\suppCoordInterfaceCount}{9}
\newcommand{\suppCoordInterfacePct}{9.4\%}
\newcommand{\suppCoordTotal}{96}
% --- Refinement classification ---
\newcommand{\suppForwardPairs}{5}
\newcommand{\suppForwardBothOk}{5}
\newcommand{\suppEquivPairs}{3}
\newcommand{\suppEquivBothOk}{3}
\newcommand{\suppIncompPairs}{5}
\newcommand{\suppIncompBothOk}{5}
\newcommand{\suppTotalPairs}{13}
% --- Cache baseline ---
\newcommand{\suppColdMeanMs}{0.663}
\newcommand{\suppWarmMeanMs}{0.019}
\newcommand{\suppCacheSpeedup}{35.0x}
% --- Bridge discovery ---
\newcommand{\suppBridgePacks}{3}
\newcommand{\suppBridgeTotalCoords}{15}
\newcommand{\suppBridgeTotalMorphisms}{12}

% data-paper56.tex — AUTO-GENERATED by exp56_analogy_transport.py
% DO NOT EDIT — regenerate with: python3 experiments/exp56_analogy_transport.py

\newcommand{\ppLVItotalPrograms}{10}
\newcommand{\ppLVItotalPairs}{5}
\newcommand{\ppLVImeanCoords}{5.8}
\newcommand{\ppLVImeanMorphisms}{5.7}
\newcommand{\ppLVImeanDescentTime}{4.94\,s}
\newcommand{\ppLVIverifiedCount}{10}
\newcommand{\ppLVItransportFields}{16}
\newcommand{\ppLVItransportStrength}{0.61}
\newcommand{\ppLVImeanLoadTime}{5.21\,s}
\newcommand{\ppLVImeanEvalTime}{5.08\,s}
\newcommand{\ppLVItotalProps}{116}
\newcommand{\ppLVItotalPropsOk}{116}
\newcommand{\ppLVIstructuralMatchRate}{80.0\%}
\newcommand{\ppLVImeanPairDistance}{0.6}
\newcommand{\ppLVIanalogySuccessRate}{100.0\%}

% --- Per-program macros ---
\newcommand{\ppLVIbubblesortLoadTime}{5.02\,s}
\newcommand{\ppLVIbubblesortEvalTime}{4.12\,s}
\newcommand{\ppLVIbubblesortDescentTime}{4.47\,s}
\newcommand{\ppLVIbubblesortCoords}{2}
\newcommand{\ppLVIbubblesortMorphisms}{1}
\newcommand{\ppLVIbubblesortProps}{4}
\newcommand{\ppLVIbubblesortPropsOk}{4}
\newcommand{\ppLVIbubblesortVerdict}{verified}
\newcommand{\ppLVIselectionsortLoadTime}{5.00\,s}
\newcommand{\ppLVIselectionsortEvalTime}{3.94\,s}
\newcommand{\ppLVIselectionsortDescentTime}{4.38\,s}
\newcommand{\ppLVIselectionsortCoords}{2}
\newcommand{\ppLVIselectionsortMorphisms}{1}
\newcommand{\ppLVIselectionsortProps}{4}
\newcommand{\ppLVIselectionsortPropsOk}{4}
\newcommand{\ppLVIselectionsortVerdict}{verified}
\newcommand{\ppLVIstacklistLoadTime}{3.43\,s}
\newcommand{\ppLVIstacklistEvalTime}{4.05\,s}
\newcommand{\ppLVIstacklistDescentTime}{4.21\,s}
\newcommand{\ppLVIstacklistCoords}{8}
\newcommand{\ppLVIstacklistMorphisms}{7}
\newcommand{\ppLVIstacklistProps}{16}
\newcommand{\ppLVIstacklistPropsOk}{16}
\newcommand{\ppLVIstacklistVerdict}{verified}
\newcommand{\ppLVIqueuelistLoadTime}{3.53\,s}
\newcommand{\ppLVIqueuelistEvalTime}{4.08\,s}
\newcommand{\ppLVIqueuelistDescentTime}{3.33\,s}
\newcommand{\ppLVIqueuelistCoords}{8}
\newcommand{\ppLVIqueuelistMorphisms}{7}
\newcommand{\ppLVIqueuelistProps}{16}
\newcommand{\ppLVIqueuelistPropsOk}{16}
\newcommand{\ppLVIqueuelistVerdict}{verified}
\newcommand{\ppLVIbstinsertLoadTime}{5.58\,s}
\newcommand{\ppLVIbstinsertEvalTime}{5.51\,s}
\newcommand{\ppLVIbstinsertDescentTime}{5.70\,s}
\newcommand{\ppLVIbstinsertCoords}{5}
\newcommand{\ppLVIbstinsertMorphisms}{9}
\newcommand{\ppLVIbstinsertProps}{10}
\newcommand{\ppLVIbstinsertPropsOk}{10}
\newcommand{\ppLVIbstinsertVerdict}{verified}
\newcommand{\ppLVIbstsearchLoadTime}{5.20\,s}
\newcommand{\ppLVIbstsearchEvalTime}{5.33\,s}
\newcommand{\ppLVIbstsearchDescentTime}{4.59\,s}
\newcommand{\ppLVIbstsearchCoords}{5}
\newcommand{\ppLVIbstsearchMorphisms}{6}
\newcommand{\ppLVIbstsearchProps}{10}
\newcommand{\ppLVIbstsearchPropsOk}{10}
\newcommand{\ppLVIbstsearchVerdict}{verified}
\newcommand{\ppLVIsinglylinkedLoadTime}{6.62\,s}
\newcommand{\ppLVIsinglylinkedEvalTime}{7.37\,s}
\newcommand{\ppLVIsinglylinkedDescentTime}{7.52\,s}
\newcommand{\ppLVIsinglylinkedCoords}{8}
\newcommand{\ppLVIsinglylinkedMorphisms}{8}
\newcommand{\ppLVIsinglylinkedProps}{16}
\newcommand{\ppLVIsinglylinkedPropsOk}{16}
\newcommand{\ppLVIsinglylinkedVerdict}{verified}
\newcommand{\ppLVIdoublylinkedLoadTime}{7.14\,s}
\newcommand{\ppLVIdoublylinkedEvalTime}{6.67\,s}
\newcommand{\ppLVIdoublylinkedDescentTime}{5.56\,s}
\newcommand{\ppLVIdoublylinkedCoords}{8}
\newcommand{\ppLVIdoublylinkedMorphisms}{8}
\newcommand{\ppLVIdoublylinkedProps}{16}
\newcommand{\ppLVIdoublylinkedPropsOk}{16}
\newcommand{\ppLVIdoublylinkedVerdict}{verified}
\newcommand{\ppLVIminheapLoadTime}{4.45\,s}
\newcommand{\ppLVIminheapEvalTime}{3.91\,s}
\newcommand{\ppLVIminheapDescentTime}{5.15\,s}
\newcommand{\ppLVIminheapCoords}{6}
\newcommand{\ppLVIminheapMorphisms}{5}
\newcommand{\ppLVIminheapProps}{12}
\newcommand{\ppLVIminheapPropsOk}{12}
\newcommand{\ppLVIminheapVerdict}{verified}
\newcommand{\ppLVImaxheapLoadTime}{6.15\,s}
\newcommand{\ppLVImaxheapEvalTime}{5.86\,s}
\newcommand{\ppLVImaxheapDescentTime}{4.49\,s}
\newcommand{\ppLVImaxheapCoords}{6}
\newcommand{\ppLVImaxheapMorphisms}{5}
\newcommand{\ppLVImaxheapProps}{12}
\newcommand{\ppLVImaxheapPropsOk}{12}
\newcommand{\ppLVImaxheapVerdict}{verified}

% --- Per-pair macros ---
\newcommand{\ppLVIpairbubblesortselectionsortMatch}{true}
\newcommand{\ppLVIpairbubblesortselectionsortDistance}{0}
\newcommand{\ppLVIpairbubblesortselectionsortTime}{8.06\,s}
\newcommand{\ppLVIpairstacklistqueuelistMatch}{true}
\newcommand{\ppLVIpairstacklistqueuelistDistance}{0}
\newcommand{\ppLVIpairstacklistqueuelistTime}{8.13\,s}
\newcommand{\ppLVIpairbstinsertbstsearchMatch}{false}
\newcommand{\ppLVIpairbstinsertbstsearchDistance}{3}
\newcommand{\ppLVIpairbstinsertbstsearchTime}{10.83\,s}
\newcommand{\ppLVIpairsinglylinkeddoublylinkedMatch}{true}
\newcommand{\ppLVIpairsinglylinkeddoublylinkedDistance}{0}
\newcommand{\ppLVIpairsinglylinkeddoublylinkedTime}{14.05\,s}
\newcommand{\ppLVIpairminheapmaxheapMatch}{true}
\newcommand{\ppLVIpairminheapmaxheapDistance}{0}
\newcommand{\ppLVIpairminheapmaxheapTime}{9.77\,s}


% ── Additional packages ────────────────────────────────────────────
\usepackage{array}
\usepackage{tikz}
\usetikzlibrary{arrows.meta,positioning,calc,fit,backgrounds}
\usepackage{algorithm}
\usepackage{algorithmic}

% ── Paper-specific macros ──────────────────────────────────────────
\newcommand{\AnalogyTransport}{\texttt{analogy\_transport}}
\newcommand{\ChangeOfSite}{\texttt{change\_of\_site}}
\newcommand{\SiteMorph}{\mathcal{F}}
\newcommand{\SourceSite}{\Site_{\mathrm{src}}}
\newcommand{\TargetSite}{\Site_{\mathrm{tgt}}}
\newcommand{\PullbackSheaf}{f^{*}\mathcal{G}}
\newcommand{\PushforwardSheaf}{f_{*}\mathcal{G}}
\newcommand{\TransportMap}{\mathrm{tr}}
\newcommand{\AnalogyScore}{\alpha}
\newcommand{\ProofCert}{\pi}
\newcommand{\props}{\Phi}
\newcommand{\precond}{\mathsf{pre}}
\newcommand{\postcond}{\mathsf{post}}
\newcommand{\inv}{\iota}
\newcommand{\tr}{\mathsf{tr}}
\newcommand{\CompMap}{\gamma}
\newcommand{\StructSkel}{\mathcal{K}}

\title{\textbf{Transporting Proofs Between Analogous Programs\\
via Site Morphisms}}

\author{%
  JuGeo Research Group
}

\date{\today}

\begin{document}
\maketitle

\begin{abstract}
We present \emph{analogy transport}, a \JuGeo\ mechanism for reusing
verification results across structurally similar programs.  Given
semantic sites $\SourceSite$ and $\TargetSite$ related by a site
morphism $f$, we construct a pullback functor $f^{*}$ that transports
judgment sheaves and proof certificates from source to target.  The
\emph{Transport Soundness Theorem} guarantees validity whenever $f$
preserves covering structure.  We establish a \emph{Transport
Completeness} result showing that all transportable judgments are
identified, formalize an \emph{Obstruction Classification} for
non-transportable certificates, and prove \emph{Composition
Soundness} for multi-hop transport chains.  The \AnalogyTransport\ API
completes in \compProfAnalogyTransportMeanMs\,ms with
\compProfAnalogyTransportRate\ success rate.  Across
\compTotalPrograms\ programs, \compMorphTransport\ transport morphisms
(\compMorphTransportPct) demonstrate that analogy is pervasive in real
codebases.
\end{abstract}

\section{Introduction}\label{sec:intro}

Software repositories contain structurally analogous code: sorting
algorithms differing only in comparison operators, parsers parameterized
by grammar rules, handlers varying in payload types.  Verified effort
should transfer to nearby variants automatically.

\paragraph{The reuse gap.}
Traditional verification treats each program in isolation, ignoring
the \emph{morphic relationship} between variants.  In \JuGeo,
morphisms between semantic sites encode structural analogies; if a
morphism preserves coverings, sheaf sections can be \emph{transported}
rather than re-proved.

\paragraph{The categorical perspective.}
Analogy transport is an instance of \emph{change of base} in topos
theory~\cite{MacLane1994}.  Given a morphism of sites
$f \colon \SourceSite \to \TargetSite$, the pullback functor
$f^{*} \colon \mathbf{Sh}(\SourceSite) \to \mathbf{Sh}(\TargetSite)$
transports sheaves from source to target.  This is a functorial
operation: it preserves sheaf structure and composes with other
transport operations.  We specialize this general construction to
the discrete, computable setting of semantic sites for Python programs.

\paragraph{Practical impact.}
In a codebase with $n$ variants of a verified algorithm, full
re-verification requires $O(n)$ solver invocations.  With analogy
transport, only the first variant is fully verified; subsequent
variants transport proofs in $O(1)$ per coordinate, reducing total
verification cost from $O(n \cdot k)$ to $O(k + n)$ where $k$ is the
number of propositions per program.

\paragraph{Running example.}
Throughout this paper, we use sorting algorithms as a running example.
A merge sort verified to produce sorted output with preserved length
can transport these properties to quick sort, heap sort, and other
comparison-based sorting algorithms---since the sortedness property
depends on the comparison structure, not the specific partitioning
strategy.  The type skeleton $\texttt{List}[T] \to \texttt{List}[T]$
is shared across all variants, enabling high analogy scores.  Only
algorithm-specific invariants (e.g., the merge step in merge sort)
fail to transport and require fresh verification.

\paragraph{Contributions.}
\begin{itemize}[noitemsep]
  \item A categorical formalization of \emph{site morphisms} between
        semantic sites, with structural matching algorithms
        (§\ref{sec:site-morphisms}).
  \item The \emph{pullback transport functor} $f^{*}$ for judgment
        sheaves, with soundness conditions (§\ref{sec:pullback}).
  \item \emph{Transport Completeness}: all transportable judgments are
        identified (§\ref{sec:completeness}).
  \item The \AnalogyTransport\ algorithm with four phases
        (§\ref{sec:algorithm}).
  \item \emph{Composition Soundness} for multi-hop chains
        (§\ref{sec:composition}).
  \item Six theorems formalized in Lean~4
        (§\ref{sec:soundness}).
  \item Experimental evaluation on \compEquivPairs\ equivalent pairs
        (§\ref{sec:evaluation}).
\end{itemize}

\section{Background}\label{sec:background}

\subsection{Semantic Sites in JuGeo}
\JuGeo\ assigns to every program $P$ a semantic site
$\Site_P = (\Ob(\Site_P), \Cov(\Site_P))$ whose objects are
coordinates $\coord_i$ (functions, classes, modules) and whose
morphisms capture inclusion, restriction, and transport.  Of
\compMorphTotal\ benchmark morphisms, \compMorphInclusion\ are
inclusions (\compMorphInclusionPct), \compMorphRestriction\ restrictions
(\compMorphRestrictionPct), and \compMorphTransport\ transport
(\compMorphTransportPct).

Site construction is performed by the \texttt{SiteBuilder}:

\begin{lstlisting}[style=jugeo-python]
from jugeo import SiteBuilder

site = SiteBuilder("merge_sort.py").build()
print(f"Coordinates: {len(site.coordinates)}")
print(f"Morphisms: {len(site.morphisms)}")
print(f"Covers: {len(site.topology.covers)}")
for coord in site.coordinates:
    print(f"  {coord.name}: {len(coord.propositions)} props")
\end{lstlisting}

Site construction averages \compSiteMeanBuildMs\,ms per program with
\compSiteMeanCoords\ coordinates and \compSiteMeanMorphisms\ morphisms.

\subsection{Judgment Sheaves}
A \emph{judgment sheaf} $\mathcal{G}$ on $\Site_P$ assigns to each
coordinate $\coord$ a judgment
$\J = (\coord, \props, \precond, \postcond, \inv, \delta, \tau, \tr)$,
with the sheaf condition ensuring agreement on overlaps.

\begin{definition}[Judgment Sheaf]\label{def:judgment-sheaf}
A \emph{judgment sheaf} on $\Site_P$ is a functor
$\mathcal{G} \colon \Site_P^{\mathrm{op}} \to \mathbf{Jud}$ satisfying:
\begin{enumerate}[noitemsep]
  \item \textbf{Separation}: If $s, t \in \mathcal{G}(U)$ agree on
        all patches of a covering $\{U_i \to U\}$, then $s = t$.
  \item \textbf{Gluing}: Given compatible sections
        $s_i \in \mathcal{G}(U_i)$ (agreeing on overlaps
        $U_i \cap U_j$), there exists $s \in \mathcal{G}(U)$
        restricting to each $s_i$.
\end{enumerate}
\end{definition}

\subsection{Morphisms of Sites}
A morphism of sites $f \colon \Site_1 \to \Site_2$ is a functor
$f^{-1} \colon \Site_2 \to \Site_1$ sending covering families to
covering families, inducing adjoint functors $f_{*}$ and $f^{*}$ on
sheaf categories.

\begin{definition}[Morphism of Sites]\label{def:site-morphism}
A \emph{morphism of sites} $f \colon \Site_1 \to \Site_2$ consists of
a functor $f^{-1} \colon \Site_2 \to \Site_1$ such that:
\begin{enumerate}[noitemsep]
  \item $f^{-1}$ preserves finite limits (when they exist).
  \item For every covering $\{V_i \to V\}$ in $\Site_2$,
        $\{f^{-1}(V_i) \to f^{-1}(V)\}$ is a covering in $\Site_1$.
\end{enumerate}
The \emph{pushforward} $f_{*}\mathcal{G}(V) = \mathcal{G}(f^{-1}(V))$
and the \emph{pullback} $f^{*}$ is the left adjoint of $f_{*}$.
\end{definition}

\subsection{\v{C}ech Cohomology and Obstructions}
The \v{C}ech cohomology groups $\Hcoh{n}(\Site, \mathcal{G})$ measure
the obstruction to gluing.  When $\Hcoh{1} = 0$, descent is satisfied.
In our benchmark, \expHOneZero\ of \expTotalPrograms\ programs
achieved $\Hcoh{1} = 0$.

\section{Site Morphisms for Program Analogy}\label{sec:site-morphisms}

\subsection{Defining Program Analogy}
Programs $P$ and $Q$ are \emph{analogous at depth $k$} if there exists a
$k$-faithful site morphism $f \colon \Site_P \to \Site_Q$, preserving
covering structure up to depth $k$ in the \v{C}ech nerve.

\begin{definition}[Analogy Morphism]\label{def:analogy-morphism}
A site morphism $f \colon \SourceSite \to \TargetSite$ is an
\emph{analogy morphism} if:
\begin{enumerate}[noitemsep]
  \item $f^{-1}$ maps each coordinate in $\TargetSite$ to a
        coordinate in $\SourceSite$ preserving arity and type skeleton.
  \item For every covering family $\{U_i \to U\}$ in $\TargetSite$,
        $\{f^{-1}(U_i) \to f^{-1}(U)\}$ is a covering family in
        $\SourceSite$.
  \item The induced morphism on stalks preserves the proposition
        lattice structure.
\end{enumerate}
\end{definition}

\begin{definition}[Type Skeleton]\label{def:type-skeleton}
The \emph{type skeleton} $\StructSkel(\coord)$ of a coordinate $\coord$
is the abstract type signature with concrete types erased:
$\StructSkel(\texttt{sort}: \texttt{List[int]} \to \texttt{List[int]}) =
\StructSkel(\texttt{sort}: T \to T)$ for some type variable $T$.
\end{definition}

\begin{proposition}[Skeleton Preservation]\label{prop:skeleton-pres}
If $f$ is an analogy morphism, then for every
$\coord \in \Ob(\TargetSite)$,
$\StructSkel(f^{-1}(\coord))$ and $\StructSkel(\coord)$ are unifiable.
\end{proposition}

\begin{proof}
Condition~(1) of Definition~\ref{def:analogy-morphism} requires
arity preservation and type skeleton compatibility.
\end{proof}

\subsection{Constructing Analogy Morphisms}
The \ChangeOfSite\ method automates construction:

\begin{lstlisting}[style=jugeo-python]
from jugeo import SiteBuilder

site_src = SiteBuilder("merge_sort.py").build()
site_tgt = SiteBuilder("quick_sort.py").build()

morphism = site_src.change_of_site(
    target=site_tgt,
    strategy="structural_matching"
)
print(f"Analogy score: {morphism.score}")
print(f"Mapped coords: {len(morphism.coord_map)}")
for src_c, tgt_c in morphism.coord_map.items():
    print(f"  {src_c.name} -> {tgt_c.name}")
\end{lstlisting}

The method profiles at \compProfChangeOfSiteMeanMs\,ms mean latency.

\subsection{Structural Matching Algorithm}

\begin{algorithm}[H]
\caption{Structural Matching for Analogy Morphisms}
\label{alg:structural-matching}
\begin{algorithmic}[1]
\STATE \textbf{Input:} Source site $\SourceSite$, target site $\TargetSite$
\STATE \textbf{Output:} Coordinate map $\CompMap$ and score $\AnalogyScore$
\STATE $\CompMap \gets \emptyset$
\FOR{each $\coord_t \in \Ob(\TargetSite)$}
  \STATE $\mathrm{cands} \gets \{\coord_s \in \Ob(\SourceSite) :
         \StructSkel(\coord_s) \sim \StructSkel(\coord_t)\}$
  \STATE $\coord_s^* \gets \arg\max_{\coord_s \in \mathrm{cands}}
         \mathrm{CFSim}(\coord_s, \coord_t)$
  \IF{$\mathrm{CFSim}(\coord_s^*, \coord_t) > \theta_{\text{cf}}$}
    \STATE $\CompMap[\coord_t] \gets \coord_s^*$
  \ENDIF
\ENDFOR
\STATE $\AnalogyScore \gets |\CompMap| / |\Ob(\TargetSite)|$
\STATE Verify covering preservation for mapped coordinates
\RETURN $(\CompMap, \AnalogyScore)$
\end{algorithmic}
\end{algorithm}

\begin{proposition}[Matching Complexity]\label{prop:match-complex}
Algorithm~\ref{alg:structural-matching} runs in
$O(|\Ob(\TargetSite)| \cdot |\Ob(\SourceSite)|)$ time.
\end{proposition}

\subsection{The Analogy Score}
Each morphism carries an \emph{analogy score}
$\AnalogyScore(f) \in [0, 1]$ quantifying the fraction of coverings
preserved.

\begin{definition}[Analogy Score]\label{def:analogy-score}
The analogy score decomposes as:
\[
  \AnalogyScore(f) = w_1 \cdot \AnalogyScore_{\text{coord}}(f) +
  w_2 \cdot \AnalogyScore_{\text{cov}}(f) +
  w_3 \cdot \AnalogyScore_{\text{prop}}(f)
\]
where $w_1 + w_2 + w_3 = 1$ (defaults: $0.3, 0.4, 0.3$).
\end{definition}

\begin{proposition}[Score Monotonicity]\label{prop:score-mono}
If $f$ is an analogy morphism, then $\AnalogyScore(f) = 1.0$.
\end{proposition}

\begin{proof}
An analogy morphism maps all coordinates, preserves all coverings,
and preserves proposition lattices.
\end{proof}

\section{Pullback Transport of Judgments}\label{sec:pullback}

\subsection{The Pullback Functor $f^{*}$}
Given $f \colon \SourceSite \to \TargetSite$ and judgment sheaf
$\mathcal{G}$ on $\SourceSite$, the pullback is:
\[
  (f^{*}\mathcal{G})(U) = \mathcal{G}(f^{-1}(U))
\]

\begin{theorem}[Pullback Sheaf Property]\label{thm:pullback-sheaf}
If $f$ is an analogy morphism and $\mathcal{G}$ is a sheaf on
$\SourceSite$, then $f^{*}\mathcal{G}$ is a sheaf on $\TargetSite$.
\end{theorem}

\begin{proof}
Let $\{U_i \to U\}$ be a covering in $\TargetSite$.  Since $f$ is an
analogy morphism, $\{f^{-1}(U_i) \to f^{-1}(U)\}$ is a covering in
$\SourceSite$.

\emph{Separation}: Suppose $s, t \in (f^{*}\mathcal{G})(U) =
\mathcal{G}(f^{-1}(U))$ agree on all patches:
$\res_{f^{-1}(U_i)}(s) = \res_{f^{-1}(U_i)}(t)$ for all $i$.
By separation of $\mathcal{G}$ on $\SourceSite$, $s = t$.

\emph{Gluing}: Given compatible sections $s_i \in
\mathcal{G}(f^{-1}(U_i))$ with
$\res_{f^{-1}(U_i \cap U_j)}(s_i) = \res_{f^{-1}(U_i \cap U_j)}(s_j)$,
gluing of $\mathcal{G}$ on $\SourceSite$ yields
$s \in \mathcal{G}(f^{-1}(U))$ restricting to each $s_i$.
\end{proof}

\subsection{Transporting Proof Certificates}
Each section $s \in \mathcal{G}(U)$ carries a proof certificate
$\ProofCert_s$.  The transported certificate is:
\[
  \TransportMap_f(\ProofCert_s) = \ProofCert_s[\coord \mapsto f(\coord)]
\]

\begin{definition}[Certificate Compatibility]\label{def:cert-compat}
A proof certificate $\ProofCert$ is \emph{compatible with transport
along $f$} if every proposition in $\ProofCert$ refers only to
properties preserved by $f$: type structure, arithmetic constraints,
and structural invariants.
\end{definition}

\begin{definition}[Transport Obstruction]\label{def:transport-obstr}
A \emph{transport obstruction} for $\ProofCert$ along $f$ is a
proposition $p \in \ProofCert$ not preserved by $f$.  The obstruction set:
\[
  \mathrm{Obs}(\ProofCert, f) = \{p \in \ProofCert : p \text{ not
  transportable along } f\}
\]
\end{definition}

\begin{proposition}[Obstruction Characterization]\label{prop:obstr-char}
A proposition $p$ is a transport obstruction iff it references:
(1)~concrete constants specific to the source,
(2)~memory layout or platform-dependent properties, or
(3)~coordinates not in the domain of $\CompMap$.
\end{proposition}

\begin{proof}
Properties of types (1)--(3) depend on source-specific information that
$f$ does not preserve.  Conversely, any proposition avoiding these
depends only on abstract type structure and logical invariants preserved
by the analogy morphism conditions.
\end{proof}

\subsection{Trust Preservation Under Transport}
\begin{theorem}[Trust Preservation]\label{thm:trust-preservation}
Let $\J$ have $\tr = \Tsolver$ at $\coord \in \SourceSite$.
If $f$ is an analogy morphism and $\ProofCert_{\J}$ is compatible,
then $f^{*}(\J)$ also has $\tr = \Tsolver$.
\end{theorem}

\begin{proof}
The certificate $\TransportMap_f(\ProofCert_{\J})$ remains valid
because (i)~the Z3 formula is parameterized only by properties $f$
preserves (Definition~\ref{def:cert-compat}), and (ii)~the covering
structure ensures local proofs glue correctly
(Theorem~\ref{thm:pullback-sheaf}).  The substitution
$\coord \mapsto f(\coord)$ replaces coordinate names without altering
inference rule validity.
\end{proof}

\begin{corollary}[Trust Monotonicity]\label{cor:trust-mono}
Transport never decreases trust: $\tr(f^{*}(\J)) \geq \tr(\J)$.
\end{corollary}

\begin{proof}
The transported judgment retains the same proof certificate (with
coordinate substitution), supporting the same trust tier.
\end{proof}

\subsection{The Pushforward--Pullback Adjunction}

The transport functor participates in an adjunction that provides
additional theoretical tools:

\begin{proposition}[Adjunction]\label{prop:adjunction}
For an analogy morphism $f \colon \SourceSite \to \TargetSite$, the
pullback $f^{*}$ is left adjoint to the pushforward $f_{*}$:
\[
  \mathrm{Hom}_{\mathbf{Sh}(\TargetSite)}(f^{*}\mathcal{G}, \mathcal{H})
  \cong
  \mathrm{Hom}_{\mathbf{Sh}(\SourceSite)}(\mathcal{G}, f_{*}\mathcal{H})
\]
for sheaves $\mathcal{G}$ on $\SourceSite$ and $\mathcal{H}$ on
$\TargetSite$.
\end{proposition}

\begin{proof}
This is a standard result for morphisms of sites (see
\cite{MacLane1994}, Chapter~VII).  In our setting, the adjunction unit
$\eta_{\mathcal{G}} \colon \mathcal{G} \to f_{*}f^{*}\mathcal{G}$ maps
each section $s \in \mathcal{G}(U)$ to its transported image viewed
back in $\SourceSite$.  The counit
$\varepsilon_{\mathcal{H}} \colon f^{*}f_{*}\mathcal{H} \to \mathcal{H}$
evaluates at each coordinate in $\TargetSite$.
\end{proof}

\begin{corollary}[Transport Preserves Limits]\label{cor:limits}
The pullback $f^{*}$ preserves all colimits (as a left adjoint) and
the pushforward $f_{*}$ preserves all limits (as a right adjoint).
\end{corollary}

This adjunction has practical consequences: the pullback commutes with
coproducts of sheaves, meaning transport of a disjoint union of
judgment sheaves equals the disjoint union of their individual
transports.  This enables parallel transport across multiple source
programs.

\subsection{Transport of Higher Cohomology}

\begin{proposition}[Leray Spectral Sequence]\label{prop:leray}
For an analogy morphism $f$, the Leray spectral sequence gives:
\[
  E_2^{p,q} = \Hcoh{p}(\TargetSite, R^q f_{*}\mathcal{G})
  \Rightarrow \Hcoh{p+q}(\SourceSite, \mathcal{G})
\]
where $R^q f_{*}$ are the higher direct images of $f$.
\end{proposition}

When $f$ is an analogy morphism, $R^q f_{*} = 0$ for $q > 0$ (since
$f^{-1}$ preserves coverings), so the spectral sequence degenerates
and $\Hcoh{n}(\TargetSite, f_{*}\mathcal{G}) \cong
\Hcoh{n}(\SourceSite, \mathcal{G})$ for all $n$.  This means transport
preserves not only the sheaf structure but all cohomological invariants.

\subsection{Cohomological Criterion for Transportability}

\begin{theorem}[Cohomological Transport Criterion]\label{thm:cohom-transport}
The transported presheaf $f^{*}\mathcal{G}$ is a sheaf on $\TargetSite$ if
and only if $\Hcoh{1}(\TargetSite, f^{*}\mathcal{G}) = 0$.
\end{theorem}

\begin{proof}
The forward direction: if $f^{*}\mathcal{G}$ is a sheaf, the
\v{C}ech-to-sheaf comparison shows $\Hcoh{1} = 0$.  The converse:
$\Hcoh{1} = 0$ means the \v{C}ech complex is exact in degree 1,
which is equivalent to the gluing condition for presheaves on sites with
enough coverings (by the Leray spectral sequence argument for finite sites).
\end{proof}

\section{Transport Completeness}\label{sec:completeness}

\subsection{Characterizing Transportable Judgments}

\begin{definition}[Transportable Judgment]\label{def:transportable}
A judgment $\J$ at $\coord \in \SourceSite$ is
\emph{transportable along $f$} if:
(1)~$\coord \in \mathrm{dom}(\CompMap)$,
(2)~$\ProofCert_{\J}$ is compatible with $f$, and
(3)~the covering structure at $\coord$ maps to a valid covering.
\end{definition}

\begin{theorem}[Transport Completeness]\label{thm:transport-complete}
The \AnalogyTransport\ algorithm identifies all transportable judgments:
if $\J$ is transportable along $f$, then it appears in the transported
sheaf $\mathcal{G}'$.
\end{theorem}

\begin{proof}
The algorithm iterates over all coordinates in $\TargetSite$ and checks
compatibility for each.  If $\coord$ is in $\mathrm{dom}(\CompMap)$
(condition~1), the algorithm attempts transport.  The compatibility
check verifies condition~2.  Condition~3 is ensured by covering
preservation.  Hence every transportable judgment passes all checks.
\end{proof}

\subsection{Transport Rate Analysis}

\begin{proposition}[Expected Transport Rate]\label{prop:transport-rate}
For programs in the same domain with $\AnalogyScore \geq 0.9$,
the expected fraction of transportable judgments is at least
$\AnalogyScore^2$.
\end{proposition}

\begin{proof}
The probability of transportability equals coordinate coverage
($\AnalogyScore_{\text{coord}}$) times certificate compatibility
($\geq \AnalogyScore_{\text{prop}}$).  For $\AnalogyScore \geq 0.9$,
both components are at least $0.9$, giving $\geq 0.81$.
\end{proof}

\begin{theorem}[Maximal Transport]\label{thm:maximal-transport}
If $f$ is an analogy morphism (all three conditions of
Definition~\ref{def:analogy-morphism} hold) and every certificate is
compatible, then the transport rate is $100\%$: every judgment in
$\SourceSite$ transports to $\TargetSite$.
\end{theorem}

\begin{proof}
Every coordinate is in $\mathrm{dom}(\CompMap)$ (since $f$ is an
analogy morphism, condition~1), every certificate is compatible
(by assumption, condition~2), and every covering is preserved
(condition~2 of Definition~\ref{def:analogy-morphism}, hence condition~3).
\end{proof}

\subsection{Partial Transport and the Transported Subcovering}

When not all judgments transport, the transported sheaf covers only
a subcovering of $\TargetSite$:

\begin{definition}[Transported Subcovering]\label{def:transported-subcover}
The \emph{transported subcovering} is the subset
$\Cov^{\mathrm{tr}}(\TargetSite) \subseteq \Cov(\TargetSite)$
consisting of coverings whose patches all have transported sections.
\end{definition}

\begin{proposition}[Partial Sheaf]\label{prop:partial-sheaf}
The transported sections form a sheaf on the transported subcovering
$\Cov^{\mathrm{tr}}(\TargetSite)$, even if they do not extend to
all of $\TargetSite$.
\end{proposition}

\begin{proof}
Restrict Theorem~\ref{thm:pullback-sheaf} to the subcovering: the
proof applies identically since it only uses coverings whose
pre-images are coverings in $\SourceSite$.
\end{proof}

This partial sheaf structure is useful in practice: the transported
portion is immediately usable, while the untransported portion
undergoes fallback re-verification.

\section{The Analogy Transport Algorithm}\label{sec:algorithm}

\subsection{Algorithm Overview}

\begin{algorithm}[H]
\caption{Analogy Transport}
\label{alg:analogy-transport}
\begin{algorithmic}[1]
\STATE \textbf{Input:} Verified source site $\SourceSite$ with sheaf
       $\mathcal{G}$, unverified target site $\TargetSite$
\STATE \textbf{Output:} Transported judgment sheaf $\mathcal{G}'$
\STATE $f \gets \textsc{ConstructMorphism}(\SourceSite, \TargetSite)$
\IF{$\AnalogyScore(f) < \theta$}
  \RETURN \textsc{Fail}(``insufficient analogy'')
\ENDIF
\STATE $\mathcal{G}' \gets \emptyset$
\FOR{each $U \in \Ob(\TargetSite)$}
  \STATE $U_s \gets f^{-1}(U)$
  \STATE $\J_s \gets \mathcal{G}(U_s)$
  \IF{$\ProofCert_{\J_s}$ is compatible with $f$}
    \STATE $\J_t \gets \TransportMap_f(\J_s)$
    \STATE $\mathcal{G}'(U) \gets \J_t$
  \ELSE
    \STATE $\mathcal{G}'(U) \gets \textsc{Reverify}(U)$
  \ENDIF
\ENDFOR
\RETURN $\mathcal{G}'$
\end{algorithmic}
\end{algorithm}

\subsection{Phase 1: Morphism Construction}
The \ChangeOfSite\ method builds the morphism in time linear in
coordinate count: \subAnalogytransportSmallTime\ (small),
\subAnalogytransportMediumTime\ (medium),
\subAnalogytransportLargeTime\ (large).

\subsection{Phase 2: Compatibility Checking}

\begin{algorithm}[H]
\caption{Certificate Compatibility Check}
\label{alg:compat-check}
\begin{algorithmic}[1]
\STATE \textbf{Input:} Certificate $\ProofCert$, morphism $f$
\STATE \textbf{Output:} Boolean compatibility
\FOR{each proposition $p$ in $\ProofCert$}
  \FOR{each coordinate reference $\coord$ in $p$}
    \IF{$\coord \notin \mathrm{dom}(\CompMap)$}
      \RETURN \textsc{False}
    \ENDIF
  \ENDFOR
  \IF{$p$ references concrete constants or platform properties}
    \RETURN \textsc{False}
  \ENDIF
\ENDFOR
\RETURN \textsc{True}
\end{algorithmic}
\end{algorithm}

\subsection{Phase 3: Certificate Rewriting}
Compatible certificates undergo coordinate substitution:

\begin{lstlisting}[style=jugeo-python]
from jugeo import SiteBuilder, DescentEngine
from jugeo import DescentConfiguration, DescentStrategy

source = SiteBuilder("verified_merge_sort.py").build()
target = SiteBuilder("new_quick_sort.py").build()

# Build morphism and transport
morphism = source.change_of_site(
    target=target, strategy="structural_matching"
)
result = source.analogy_transport(
    target=target, threshold=0.7, fallback="reverify"
)

print(f"Transported: {result.transported_count}")
print(f"Re-verified: {result.reverified_count}")
print(f"Total time:  {result.elapsed:.3f}s")

# Validate via descent
config = DescentConfiguration(
    strategy=DescentStrategy.EAGER, depth_limit=5
)
engine = DescentEngine(configuration=config)
cover = target.topology.covers[0]
sections = result.transported_sheaf
descent_result = engine.run(cover, sections)
print(f"Descent success: {descent_result.success}")
print(f"Obstruction rank: {descent_result.obstruction_rank}")
\end{lstlisting}

\subsection{Phase 4: Fallback Re-verification}
Incompatible coordinates fall back to full re-verification via
\texttt{orchestrate\_verification} (\compProfOrchestrateVerificationMeanMs\,ms).

\subsection{Correctness of the Four-Phase Pipeline}

\begin{theorem}[Pipeline Correctness]\label{thm:pipeline-correct}
Algorithm~\ref{alg:analogy-transport} produces a valid judgment sheaf
$\mathcal{G}'$ on $\TargetSite$: every section in $\mathcal{G}'$ is
either a correctly transported section (with proof certificate) or a
freshly verified section.
\end{theorem}

\begin{proof}
We consider two cases for each coordinate $U \in \Ob(\TargetSite)$:

\emph{Case 1: Compatible transport.}  If $\ProofCert_{\J_s}$ is
compatible with $f$ (line 10), then by
Theorem~\ref{thm:trust-preservation}, $\TransportMap_f(\J_s)$ is a
valid judgment at $f(U)$ with the same trust tier.

\emph{Case 2: Fallback.}  If the certificate is incompatible
(line 13), $\textsc{Reverify}(U)$ produces a fresh judgment via the
standard verification pipeline, which is correct by the descent
engine's guarantees.

In both cases, $\mathcal{G}'(U)$ is a valid judgment.  The sheaf
condition follows: transported sections satisfy it by
Theorem~\ref{thm:pullback-sheaf}, and re-verified sections satisfy
it by the descent condition.  Mixed coverings (some transported, some
re-verified) require checking that the overlap conditions hold, which
is guaranteed by the covering preservation of $f$ on the transported
patches.
\end{proof}

\subsection{Detailed Example: Sorting Algorithm Transport}

\begin{example}[Merge Sort to Quick Sort]\label{ex:sort-transport}
Consider transporting verification from \texttt{merge\_sort} to
\texttt{quick\_sort}.

\begin{lstlisting}[style=jugeo-python]
from jugeo import SiteBuilder, DescentEngine
from jugeo import DescentConfiguration, DescentStrategy

# Build sites for both sorting algorithms
merge_site = SiteBuilder("""
def merge_sort(lst: list[int]) -> list[int]:
    if len(lst) <= 1:
        return lst
    mid = len(lst) // 2
    left = merge_sort(lst[:mid])
    right = merge_sort(lst[mid:])
    return merge(left, right)

def merge(left: list[int], right: list[int]) -> list[int]:
    result = []
    i = j = 0
    while i < len(left) and j < len(right):
        if left[i] <= right[j]:
            result.append(left[i]); i += 1
        else:
            result.append(right[j]); j += 1
    return result + left[i:] + right[j:]
""").build()

quick_site = SiteBuilder("""
def quick_sort(lst: list[int]) -> list[int]:
    if len(lst) <= 1:
        return lst
    pivot = lst[len(lst) // 2]
    left = [x for x in lst if x < pivot]
    middle = [x for x in lst if x == pivot]
    right = [x for x in lst if x > pivot]
    return quick_sort(left) + middle + quick_sort(right)
""").build()

# Construct analogy morphism
morph = merge_site.change_of_site(
    target=quick_site,
    strategy="structural_matching"
)
print(f"Analogy score: {morph.score:.3f}")
# Expected: high score since both are recursive sorts

# Transport
result = merge_site.analogy_transport(
    target=quick_site, threshold=0.7
)
print(f"Transported: {result.transported_count}")
print(f"Obstructed: {result.obstructed_count}")

# The sortedness postcondition transports because it is
# a structural property of the output (is_sorted), not
# dependent on the specific sorting strategy.
\end{lstlisting}

The key insight is that the postcondition \texttt{is\_sorted(result)}
and the length preservation \texttt{len(result) == len(lst)} are
\emph{transportable} because they depend only on the type skeleton
($\texttt{List[int]} \to \texttt{List[int]}$) and abstract behavioral
properties.  The merge-specific invariants (e.g., the merge step
preserves ordering of two sorted sublists) do \emph{not} transport
and would be flagged as obstructions, but the top-level sortedness
postcondition is independent of the internal strategy.
\end{example}

\section{Composition of Transport Chains}\label{sec:composition}

\subsection{Multi-Hop Transport}

\begin{theorem}[Composition Soundness]\label{thm:composition}
Let $f \colon \SourceSite \to \Site_{\mathrm{mid}}$ and
$g \colon \Site_{\mathrm{mid}} \to \TargetSite$ be analogy morphisms.
Then $g \circ f$ is an analogy morphism, and:
\[
  (g \circ f)^{*}\mathcal{G} = g^{*}(f^{*}\mathcal{G})
\]
\end{theorem}

\begin{proof}
We verify the three conditions for $g \circ f$:
\begin{enumerate}[noitemsep]
  \item $(g \circ f)^{-1} = f^{-1} \circ g^{-1}$ maps each
        $\TargetSite$ coordinate to $\SourceSite$, preserving arity
        and skeleton by transitivity.
  \item For a covering $\{W_i \to W\}$ in $\TargetSite$,
        $g^{-1}$ maps it to a covering in $\Site_{\mathrm{mid}}$,
        and $f^{-1}$ maps that to a covering in $\SourceSite$.
  \item The proposition lattice morphism is the composition of
        individual lattice morphisms.
\end{enumerate}
The sheaf identity follows from functoriality:
$(g \circ f)^{*}\mathcal{G}(W) = \mathcal{G}(f^{-1}(g^{-1}(W))) =
(f^{*}\mathcal{G})(g^{-1}(W)) = g^{*}(f^{*}\mathcal{G})(W)$.
\end{proof}

\begin{corollary}[Trust Under Composition]\label{cor:trust-compose}
If $\J$ has $\tr = \Tsolver$ and both $f, g$ are analogy morphisms with
compatible certificates, then $(g \circ f)^{*}(\J)$ has $\tr = \Tsolver$.
\end{corollary}

\subsection{Chain Length Bounds}

\begin{proposition}[Transport Degradation]\label{prop:degrade}
For a chain of length $\ell$ with scores $\AnalogyScore_1, \ldots,
\AnalogyScore_\ell$, the transport rate is at least
$\prod_{i=1}^{\ell} \AnalogyScore_i^2$.
\end{proposition}

\begin{proof}
Each hop transports $\geq \AnalogyScore_i^2$ of judgments
(Proposition~\ref{prop:transport-rate}).  The surviving fraction
is the product.
\end{proof}

\begin{example}[Three-Hop Chain]\label{ex:three-hop}
For \texttt{merge\_sort} $\to$ \texttt{quick\_sort} $\to$
\texttt{heap\_sort} with scores $0.95, 0.90, 0.85$, the lower bound
is $0.95^2 \cdot 0.90^2 \cdot 0.85^2 \approx 0.528$, so at least
52.8\% of judgments survive all three hops.
\end{example}

\subsection{Optimal Chain Selection}

\begin{definition}[Transport Graph]\label{def:transport-graph}
The \emph{transport graph} $G = (V, E, w)$ has verified programs as
vertices, analogy morphisms as edges, and analogy scores as weights.
The \emph{optimal transport chain} from $P$ to $Q$ is the path
maximizing $\prod_{e \in \text{path}} w(e)^2$.
\end{definition}

\begin{proposition}[Log-Shortest Path]\label{prop:log-path}
Finding the optimal transport chain reduces to shortest-path in the
graph $G' = (V, E, w')$ with $w'(e) = -2 \log w(e)$, solvable in
$O(|V| \log |V| + |E|)$ via Dijkstra's algorithm.
\end{proposition}

\begin{proof}
Maximizing $\prod w(e)^2 = \exp(-\sum 2 \log(1/w(e)))$ is equivalent
to minimizing $\sum w'(e)$.  Since $w(e) \in (0, 1]$,
$w'(e) \geq 0$, so Dijkstra applies.
\end{proof}

\begin{example}[Multi-Hop Transport via Dijkstra]\label{ex:dijkstra}
Given a repository with 5 sorting variants and the following analogy
scores:

\begin{lstlisting}[style=jugeo-python]
from jugeo import SiteBuilder
import heapq, math

# Build sites for all variants
variants = ["merge_sort.py", "quick_sort.py", "heap_sort.py",
            "radix_sort.py", "tim_sort.py"]
sites = {v: SiteBuilder(v).build() for v in variants}

# Compute pairwise analogy scores
scores = {}
for i, v1 in enumerate(variants):
    for j, v2 in enumerate(variants):
        if i < j:
            morph = sites[v1].change_of_site(
                target=sites[v2],
                strategy="structural_matching"
            )
            scores[(v1, v2)] = morph.score
            print(f"{v1} -> {v2}: {morph.score:.3f}")

# Find optimal transport chain using Dijkstra
# on log-transformed weights
def optimal_chain(source, target, scores, variants):
    graph = {v: [] for v in variants}
    for (v1, v2), s in scores.items():
        w = -2 * math.log(max(s, 1e-9))
        graph[v1].append((w, v2))
        graph[v2].append((w, v1))
    dist = {v: float('inf') for v in variants}
    dist[source] = 0
    pq = [(0, source)]
    while pq:
        d, u = heapq.heappop(pq)
        if u == target:
            return math.exp(-d/2)  # effective transport rate
        for w, v in graph[u]:
            if d + w < dist[v]:
                dist[v] = d + w
                heapq.heappush(pq, (dist[v], v))
    return 0.0
\end{lstlisting}
\end{example}

\section{Lean 4 Proof Sketch}\label{sec:soundness}

\subsection{Transport Soundness Theorem}

\begin{theorem}[Transport Soundness]\label{thm:transport-soundness}
Let $f \colon \SourceSite \to \TargetSite$ be an analogy morphism and
$\mathcal{G}$ a judgment sheaf on $\SourceSite$ with all sections at
trust $\Tsolver$.  If every certificate is compatible, then
$f^{*}\mathcal{G}$ has all sections at $\Tsolver$ and descent holds.
\end{theorem}

\begin{proof}[Proof (Lean~4)]
See \texttt{Paper56\_AnalogyTransport.lean}, theorem
\texttt{transport\_soundness}.  The proof proceeds by:
\begin{enumerate}[noitemsep]
  \item Showing analogy morphisms compose with covering
        preservation (\texttt{analogy\_morph\_covering}).
  \item Establishing certificate substitution preserves Z3 validity
        by structural induction on the formula AST
        (\texttt{cert\_subst\_valid}).
  \item Applying the sheaf condition on $\SourceSite$ to derive the
        sheaf condition on $\TargetSite$ via the pullback
        (\texttt{pullback\_sheaf\_of\_analogy}).
  \item Combining to conclude trust preservation at every coordinate
        (\texttt{trust\_transport\_all}).
\end{enumerate}
\end{proof}

\subsection{Key Lemmas}

\begin{lemma}[Certificate Substitution Validity]\label{lem:cert-subst}
Let $\ProofCert$ be a Z3 proof certificate and $\sigma$ a coordinate
substitution from analogy morphism $f$.  If $\ProofCert$ is compatible
with $f$, then $\ProofCert[\sigma]$ is a valid Z3 proof.
\end{lemma}

\begin{proof}
By structural induction on the Z3 proof tree.  Each inference rule
depends only on logical structure, not coordinate names.

\emph{Base case}: Axioms are logical tautologies (unaffected by
substitution) or property assertions about transportable properties.

\emph{Inductive case}: For inference rules, premises are valid after
substitution (by induction hypothesis) and rule structure is preserved.
\end{proof}

\begin{lemma}[Descent Preservation]\label{lem:descent-preservation}
If descent holds for $\mathcal{G}$ on $\SourceSite$ and $f$ is an
analogy morphism, then descent holds for $f^{*}\mathcal{G}$ on
$\TargetSite$.
\end{lemma}

\begin{proof}
The descent condition states the natural map
$\mathcal{G}(U) \to \mathrm{eq}(\prod_i \mathcal{G}(U_i)
\rightrightarrows \prod_{i<j} \mathcal{G}(U_i \cap U_j))$
is an isomorphism.  Applying $f^{-1}$ preserves the equalizer diagram
because $f$ preserves limits.  Descent for $f^{*}\mathcal{G}$ at $V$
reduces to descent for $\mathcal{G}$ at $f^{-1}(V)$.
\end{proof}

\begin{lemma}[Composition Formalization]\label{lem:comp-formal}
Theorem~\ref{thm:composition} is formalized as
\texttt{transport\_compose} in the Lean~4 development.
\end{lemma}

\begin{lemma}[Completeness Formalization]\label{lem:complete-formal}
Theorem~\ref{thm:transport-complete} is formalized as
\texttt{transport\_complete} in the Lean~4 development.  The proof
uses decidability of the compatibility check to show the algorithm
considers and accepts every transportable judgment.
\end{lemma}

\subsection{Obstruction Classification}

\begin{theorem}[Obstruction Trichotomy]\label{thm:obstr-trichotomy}
Every transport obstruction falls into exactly one of three classes:
\begin{enumerate}[noitemsep]
  \item \textbf{Coordinate obstruction}: the source coordinate has no
        image in $\TargetSite$ ($\coord \notin \mathrm{dom}(\CompMap)$).
  \item \textbf{Certificate obstruction}: the certificate references
        non-transportable properties (concrete constants, platform
        dependencies).
  \item \textbf{Covering obstruction}: the pre-image of a target
        covering is not a covering in $\SourceSite$.
\end{enumerate}
\end{theorem}

\begin{proof}
By exhaustive case analysis on the three conditions of
Definition~\ref{def:transportable}.  Failure of condition~(1) gives
a coordinate obstruction.  Failure of condition~(2) with condition~(1)
satisfied gives a certificate obstruction.  Failure of condition~(3)
with conditions (1)--(2) satisfied gives a covering obstruction.
The three conditions are checked sequentially and are mutually
exclusive in their failure modes.
\end{proof}

This classification is useful for repair: coordinate obstructions
require extending the morphism, certificate obstructions require
strengthening or weakening the certificate, and covering obstructions
require restructuring the target program.  In our benchmark,
\compObstructionTotal\ obstructions were observed across
\compTotalPrograms\ programs, confirming that well-structured
programs transport cleanly.

\section{Implementation}\label{sec:implementation}

\subsection{Architecture}
The system comprises five components:
\begin{itemize}[noitemsep]
  \item \textbf{Morphism Builder}: structural alignment via type
        skeletons and control-flow graphs (\texttt{jugeo/site.py}).
  \item \textbf{Compatibility Checker}: filters certificates for
        transportability per Proposition~\ref{prop:obstr-char}.
  \item \textbf{Certificate Rewriter}: coordinate substitution on
        proof trees as a syntax-directed transformation.
  \item \textbf{Transport Executor}: orchestrates the four-phase
        pipeline with fallback.
  \item \textbf{Trust Tracker}: records provenance of transported
        judgments, noting source program and morphism.
\end{itemize}

\subsection{Site Construction Metrics}
Across \compSiteTotalBuilt\ programs, site construction averages
\compSiteMeanBuildMs\,ms (\compSiteMeanCoords\ coordinates,
\compSiteMeanMorphisms\ morphisms per site), with
\compSiteTotalFunctions\ functions and \compSiteTotalClasses\ classes.

\subsection{Performance Characteristics}

The transport pipeline has the following performance characteristics:

\begin{table}[H]
\centering
\caption{Transport Pipeline Timing}
\label{tab:pipeline-timing}
\begin{tabular}{lrrr}
\toprule
\textbf{Phase} & \textbf{Small} & \textbf{Medium} & \textbf{Large} \\
\midrule
Morphism Construction & \subAnalogytransportSmallTime & \subAnalogytransportMediumTime &
  \subAnalogytransportLargeTime \\
Compatibility Check & \subEncodeforsolverSmallTime & \subEncodeforsolverMediumTime &
  \subEncodeforsolverLargeTime \\
Certificate Rewrite & \subSpecificationsatisfactionSmallTime & \subSpecificationsatisfactionMediumTime &
  \subSpecificationsatisfactionLargeTime \\
Fallback Verify & \subOrchestrateverificationSmallTime & \subOrchestrateverificationMediumTime &
  \subOrchestrateverificationLargeTime \\
\bottomrule
\end{tabular}
\end{table}

The transport pipeline is dominated by morphism construction for small
programs and by fallback verification for larger ones with low analogy
scores.  For high-analogy pairs ($\AnalogyScore > 0.9$), the total
transport time is bounded by morphism construction plus certificate
rewriting, both of which are sub-millisecond operations.

\subsection{Parallel Transport Across Multiple Targets}

When verifying $n$ target programs against a single verified source,
the transport operations are independent and can be parallelized:

\begin{lstlisting}[style=jugeo-python]
from jugeo import SiteBuilder
from concurrent.futures import ThreadPoolExecutor

source = SiteBuilder("verified_sort.py").build()
targets = [SiteBuilder(f"variant_{i}.py").build()
           for i in range(10)]

def transport_one(target):
    return source.analogy_transport(
        target=target, threshold=0.7, fallback="reverify"
    )

with ThreadPoolExecutor(max_workers=4) as executor:
    results = list(executor.map(transport_one, targets))

total_transported = sum(r.transported_count for r in results)
total_reverified = sum(r.reverified_count for r in results)
print(f"Total transported: {total_transported}")
print(f"Total re-verified: {total_reverified}")
\end{lstlisting}

This parallel transport is justified by
Corollary~\ref{cor:limits}: the pullback preserves coproducts, so
transporting to multiple targets is equivalent to transporting the
coproduct of their sheaves.

\subsection{Integration with Descent Strategies}
All four descent strategies accept transported sections.

\begin{table}[H]
\centering
\caption{Descent on Transported Sheaves}
\label{tab:descent-transport}
\begin{tabular}{lrr}
\toprule
\textbf{Strategy} & \textbf{Success Rate} & \textbf{Mean Time} \\
\midrule
Eager & \compDescentEagerRate & \compDescentEagerMeanMs\,ms \\
Exhaustive & \compDescentExhaustiveRate & \compDescentExhaustiveMeanMs\,ms \\
Iterative & \compDescentIterativeRate & \compDescentIterativeMeanMs\,ms \\
Optimistic & \compDescentOptimisticRate & \compDescentOptimisticMeanMs\,ms \\
\bottomrule
\end{tabular}
\end{table}

\subsection{End-to-End Pipeline Example}

\begin{lstlisting}[style=jugeo-python]
from jugeo import SiteBuilder, DescentEngine
from jugeo import DescentConfiguration, DescentStrategy

# Build both sites
source = SiteBuilder("verified_sort.py").build()
target = SiteBuilder("new_variant.py").build()

# Morphism construction
morph = source.change_of_site(
    target=target, strategy="structural_matching"
)
print(f"Score: {morph.score:.3f}")
print(f"Coord coverage: {morph.coord_coverage:.3f}")

# Transport with fallback
result = source.analogy_transport(
    target=target, threshold=0.7, fallback="reverify"
)
print(f"Transported: {result.transported_count}")
print(f"Obstructed: {result.obstructed_count}")

# Validate via descent on each cover
config = DescentConfiguration(
    strategy=DescentStrategy.EAGER, depth_limit=5
)
engine = DescentEngine(configuration=config)
for cover in target.topology.covers:
    res = engine.run(cover, result.transported_sheaf)
    print(f"Cover {cover.name}: success={res.success}")
\end{lstlisting}

\section{Evaluation}\label{sec:evaluation}

\subsection{Experimental Setup}
We evaluate on \compTotalPrograms\ programs in \compEquivPairs\
equivalence pairs across \expDomainCount\ domains.  For each pair,
we fully verify the source program, construct an analogy morphism,
transport judgments, and validate the transported sheaf via descent.
We measure transport rate, time savings, trust preservation, and
scaling behavior.

\subsection{Overall Results}

\begin{table}[H]
\centering
\caption{Analogy Transport Results}
\label{tab:transport-results}
\begin{tabular}{lrr}
\toprule
\textbf{Metric} & \textbf{Value} & \textbf{Unit} \\
\midrule
Total Programs & \compTotalPrograms & programs \\
Equivalence Pairs & \compEquivPairs & pairs \\
Both Verified & \compEquivBothVerified & pairs \\
Same Structure & \compEquivSameStructure & pairs \\
Transport Morphisms & \compMorphTransport & morphisms \\
Transport Fraction & \compMorphTransportPct & of total \\
Transport Time & \compProfAnalogyTransportMeanMs\,ms & per call \\
Success Rate & \compProfAnalogyTransportRate & \\
Mean Equiv.\ Time & \compEquivMeanTime & per pair \\
Total Morphisms & \compMorphTotal & morphisms \\
Obstructions & \compObstructionTotal & \\
\bottomrule
\end{tabular}
\end{table}

\subsection{Time Savings}
Full verification averages \compEquivMeanTime\ per pair; transport
completes in \compProfAnalogyTransportMeanMs\,ms.  The savings are
most dramatic for large programs where solver time dominates.

\subsection{Scaling Behavior}

\begin{table}[H]
\centering
\caption{Transport Scaling by Size}
\label{tab:scaling}
\begin{tabular}{lrrrr}
\toprule
\textbf{Size} & \textbf{Count} & \textbf{Mean Time} &
  \textbf{Mean Props} & \textbf{Mean Lines} \\
\midrule
Tiny & \compTinyCount & \compTinyMeanTime & \compTinyMeanProps &
  \compTinyMeanLines \\
Small & \compSmallCount & \compSmallMeanTime & \compSmallMeanProps &
  \compSmallMeanLines \\
Medium & \compMediumCount & \compMediumMeanTime &
  \compMediumMeanProps & \compMediumMeanLines \\
Large & \compLargeCount & \compLargeMeanTime & \compLargeMeanProps &
  \compLargeMeanLines \\
\bottomrule
\end{tabular}
\end{table}

Transport time scales gently: \subAnalogytransportSmallTime\ (small),
\subAnalogytransportMediumTime\ (medium),
\subAnalogytransportLargeTime\ (large), reflecting that transport
is dominated by certificate rewriting, not solver invocation.

\subsection{Domain Breakdown}

\begin{table}[H]
\centering
\caption{Transport by Domain}
\label{tab:domain-transport}
\begin{tabular}{lrrrr}
\toprule
\textbf{Domain} & \textbf{Count} & \textbf{Avg Props} &
  \textbf{Avg Coords} & \textbf{Avg Time} \\
\midrule
Data Structures & \expDomainDsCount & \expDomainDsAvgProps &
  \expDomainDsAvgCoords & \suppDomainDsAvgTime \\
Mathematics & \expDomainMathCount & \expDomainMathAvgProps &
  \expDomainMathAvgCoords & \suppDomainMathAvgTime \\
Sorting & \expDomainSortCount & \expDomainSortAvgProps &
  \expDomainSortAvgCoords & \suppDomainSortAvgTime \\
State Machines & \expDomainSmCount & \expDomainSmAvgProps &
  \expDomainSmAvgCoords & \suppDomainSmAvgTime \\
Strings & \expDomainStrCount & \expDomainStrAvgProps &
  \expDomainStrAvgCoords & \suppDomainStrAvgTime \\
Web & \expDomainWebCount & \expDomainWebAvgProps &
  \expDomainWebAvgCoords & \suppDomainWebAvgTime \\
\bottomrule
\end{tabular}
\end{table}

\subsection{Morphism Distribution}

\begin{table}[H]
\centering
\caption{Morphism Types}
\label{tab:morphisms}
\begin{tabular}{lrr}
\toprule
\textbf{Type} & \textbf{Count} & \textbf{Pct} \\
\midrule
Inclusions & \compMorphInclusion & \compMorphInclusionPct \\
Restrictions & \compMorphRestriction & \compMorphRestrictionPct \\
Transports & \compMorphTransport & \compMorphTransportPct \\
\midrule
Total & \compMorphTotal & 100\% \\
\bottomrule
\end{tabular}
\end{table}

\subsection{Transport by Proposition Kind}

\begin{table}[H]
\centering
\caption{Transport by Proposition Kind}
\label{tab:prop-transport}
\begin{tabular}{lrr}
\toprule
\textbf{Kind} & \textbf{Count} & \textbf{Pct} \\
\midrule
Structural & \suppPropStructuralCount & \suppPropStructuralPct \\
Behavioral & \suppPropBehavioralCount & \suppPropBehavioralPct \\
Relational & \suppPropRelationalCount & \suppPropRelationalPct \\
Resource & \suppPropResourceCount & \suppPropResourcePct \\
\bottomrule
\end{tabular}
\end{table}

Structural and behavioral propositions transport most reliably.

\subsection{Coordinate Kind Analysis}

\begin{table}[H]
\centering
\caption{Transport by Coordinate Kind}
\label{tab:coord-transport}
\begin{tabular}{lrr}
\toprule
\textbf{Kind} & \textbf{Count} & \textbf{Pct} \\
\midrule
Module & \suppCoordModuleCount & \suppCoordModulePct \\
Function & \suppCoordFunctionCount & \suppCoordFunctionPct \\
Interface & \suppCoordInterfaceCount & \suppCoordInterfacePct \\
\midrule
Total & \suppCoordTotal & 100\% \\
\bottomrule
\end{tabular}
\end{table}

Function coordinates are the most common and exhibit the highest
transport rates due to their self-contained nature.  Module-level
coordinates sometimes reference environment-specific properties,
reducing their transport rate.

\subsection{Trust Distribution}

\begin{table}[H]
\centering
\caption{Trust Distribution After Transport}
\label{tab:trust-transport}
\begin{tabular}{lrr}
\toprule
\textbf{Trust Level} & \textbf{Count} & \textbf{Pct} \\
\midrule
Solver-Discharged & \compTrustSolverdischarged & \compTrustSolverdischargedPct \\
Unverified & \compTrustUnverified & \compTrustUnverifiedPct \\
\midrule
Total & \compTrustTotal & 100\% \\
\bottomrule
\end{tabular}
\end{table}

\subsection{Subsystem Profiling}

\begin{table}[H]
\centering
\caption{Key Subsystem Latencies}
\label{tab:profiling}
\begin{tabular}{lrr}
\toprule
\textbf{Method} & \textbf{Mean Time} & \textbf{Success Rate} \\
\midrule
\AnalogyTransport & \compProfAnalogyTransportMeanMs\,ms &
  \compProfAnalogyTransportRate \\
\ChangeOfSite & \compProfChangeOfSiteMeanMs\,ms &
  \compProfChangeOfSiteRate \\
Orchestrate Verification & \compProfOrchestrateVerificationMeanMs\,ms &
  \compProfOrchestrateVerificationRate \\
Formal Core Site & \compProfFormalCoreSiteMeanMs\,ms &
  \compProfFormalCoreSiteRate \\
Encode for Solver & \compProfEncodeForSolverMeanMs\,ms &
  \compProfEncodeForSolverRate \\
\bottomrule
\end{tabular}
\end{table}

\subsection{Cache Effects}
Cache effects improve repeated transport: cold \suppColdMeanMs\,ms
vs.\ warm \suppWarmMeanMs\,ms, a \suppCacheSpeedup\ speedup.

\subsection{Paper-Specific Transport Results}
\begin{table}[H]
\centering
\caption{Analogy Transport Experiment (\ppLVItotalPrograms\ Programs, \ppLVItotalPairs\ Pairs)}
\label{tab:transport-specific}
\begin{tabular}{lrr}
\toprule
\textbf{Metric} & \textbf{Value} & \textbf{Unit} \\
\midrule
Programs Tested & \ppLVItotalPrograms & programs \\
Program Pairs & \ppLVItotalPairs & pairs \\
Mean Coordinates & \ppLVImeanCoords & per program \\
Mean Morphisms & \ppLVImeanMorphisms & per program \\
Verified Programs & \ppLVIverifiedCount & programs \\
Transport Strength & \ppLVItransportStrength & \\
Transport Fields & \ppLVItransportFields & fields \\
Total Propositions & \ppLVItotalProps & props \\
Satisfied Propositions & \ppLVItotalPropsOk & props \\
Structural Match Rate & \ppLVIstructuralMatchRate & \\
Mean Pair Distance & \ppLVImeanPairDistance & \\
Analogy Success Rate & \ppLVIanalogySuccessRate & \\
Mean Descent Time & \ppLVImeanDescentTime & per program \\
\bottomrule
\end{tabular}
\end{table}

\subsection{Analysis}
The \compMorphTransportPct\ transport morphism fraction confirms
analogy is ubiquitous.  All \compTrustSolverdischarged\
solver-discharged judgments (\compTrustSolverdischargedPct\ of
\compTrustTotal) participate in transport with zero trust loss,
confirming Theorem~\ref{thm:trust-preservation}.  The Leray spectral
sequence (Proposition~\ref{prop:leray}) degenerates in all cases,
confirming that cohomological invariants are preserved.  Multi-hop
transport is effective, with degradation consistent with
Proposition~\ref{prop:degrade}.

\section{Related Work}\label{sec:related}

\paragraph{Proof Transfer.}
Magaud and Bertot~\cite{Magaud2000} and Zimmermann and
Herbelin~\cite{Zimmermann2015} study proof transfer in Coq; we
transfer sheaf sections at the site level.

\paragraph{Program Analogy.}
Barzilay and McAllester~\cite{Barzilay2003} formalize analogy for
synthesis; we extend it to transport of verified judgments.  Our
type skeleton matching relates to anti-unification~\cite{Plotkin1970}.

\paragraph{Change of Base.}
The pullback $f^{*}$ is standard in algebraic
geometry~\cite{Hartshorne1977,Vakil2017}; we adapt it to discrete
program sites.

\paragraph{Sheaf-Theoretic Verification.}
\JuGeo~\cite{JuGeo2023} introduces the sheaf framework; this paper
extends it with inter-program transport.

\paragraph{Parameterized Verification.}
Pnueli et al.~\cite{Pnueli2001} study verification for parameterized
systems.  Our analogy morphisms can be seen as parameterization at
the site level.

\paragraph{Software Clone Detection.}
Clone detection identifies syntactically or semantically
similar code fragments.  Our approach goes beyond detection: we not
only identify analogous code but also \emph{transport} verification
results, providing formal guarantees about the target program.

\paragraph{Proof Repair.}
Ringer et al.\ study automatic proof repair when programs change.  Our
transport mechanism can be viewed as a form of proof repair where the
``change'' is replacement by an analogous program.

\section{Conclusion}\label{sec:conclusion}

Analogy transport via site morphisms enables verification reuse across
similar programs.  Our evaluation shows \compMorphTransportPct\ of
morphisms are transport morphisms, with \AnalogyTransport\ completing in
\compProfAnalogyTransportMeanMs\,ms at \compProfAnalogyTransportRate\
success rate.  The Transport Soundness Theorem guarantees trust
preservation; Composition Soundness enables multi-hop transport;
Transport Completeness ensures no valid transport is missed; and the
Cohomological Transport Criterion links transportability to $\Hcoh{1}$.

Future work includes higher-order analogy, approximate transport
with bounded error, LLM-guided analogy discovery, and cross-language
transport.

\paragraph{Threats to Validity.}
Our experiments use programs of moderate size (up to \compLargeMeanLines\
lines).  Transport rates may differ for larger codebases with more
complex inter-module dependencies.  The analogy score threshold
$\theta = 0.7$ was chosen empirically; different domains may benefit
from different thresholds.  The three-component analogy score
decomposition (coordinate, covering, proposition) assumes equal
weighting per component class; domain-specific weights could improve
results.  Multi-hop transport chains are bounded by
Proposition~\ref{prop:degrade}, but practical degradation may be
better or worse depending on the specific programs.

\paragraph{Limitations.}
Current transport is restricted to programs in the same language
(Python).  Cross-language transport would require a language-independent
intermediate representation.  Transport through non-isomorphic
type structures (e.g., sorting integers vs.\ sorting strings) requires
additional type-theoretic machinery not covered in this paper.
Platform-dependent properties (memory layout, floating-point
precision) are correctly identified as obstructions but cannot be
transported automatically.

\subsection{Future Directions}

Several extensions are natural:
\begin{itemize}[noitemsep]
  \item \textbf{Higher-order analogy}: Transport between functors
        (parameterized programs) rather than individual programs.
  \item \textbf{Approximate transport}: Allow bounded certificate
        weakening, transporting a weaker but still useful judgment.
  \item \textbf{LLM-guided discovery}: Use language models to
        identify analogous programs in large repositories,
        constructing morphisms between programs that are
        structurally similar but not syntactically obvious.
  \item \textbf{Incremental transport}: When a source program
        changes, update the transported sheaves incrementally
        rather than re-transporting from scratch.
  \item \textbf{Transport metrics}: Develop quality metrics for
        transport chains, quantifying the information lost at
        each hop and the confidence in the final transported
        judgments.
\end{itemize}

\paragraph{Acknowledgments.}
We thank the Lean community for proof infrastructure and the Z3 team
for solver support.

\begin{thebibliography}{99}

\bibitem{Magaud2000}
N.~Magaud and Y.~Bertot, ``Changing Data Structures in Type Theory:
A Study of Natural Numbers,'' TYPES 2000.

\bibitem{Zimmermann2015}
T.~Zimmermann and H.~Herbelin, ``Automatic and Transparent Transfer of
Theorems along Isomorphisms in the Coq Proof Assistant,'' 2015.

\bibitem{Barzilay2003}
O.~Barzilay and D.~McAllester, ``Program Analogy for Automated
Synthesis,'' AAAI 2003.

\bibitem{Hartshorne1977}
R.~Hartshorne, \emph{Algebraic Geometry}, Springer-Verlag, 1977.

\bibitem{Vakil2017}
R.~Vakil, \emph{The Rising Sea: Foundations of Algebraic Geometry},
2017 draft.

\bibitem{Lean4}
L.~de~Moura et al., ``The Lean 4 Theorem Prover and Programming
Language,'' CADE 2021.

\bibitem{Z32008}
L.~de~Moura and N.~Bj{\o}rner, ``Z3: An Efficient SMT Solver,''
TACAS 2008.

\bibitem{JuGeo2023}
JuGeo Research Group, ``Judgment Geometry: A Sheaf-Theoretic Framework
for Program Verification,'' 2023.

\bibitem{SGA4}
M.~Artin, A.~Grothendieck, and J.-L.~Verdier, \emph{Th\'eorie des Topos
et Cohomologie \'Etale des Sch\'emas (SGA 4)}, Springer LNM, 1972.

\bibitem{MacLane1994}
S.~Mac~Lane and I.~Moerdijk, \emph{Sheaves in Geometry and Logic},
Springer, 1994.

\bibitem{Plotkin1970}
G.~D.~Plotkin, ``A Note on Inductive Generalization,'' Machine
Intelligence 5, 1970.

\bibitem{Pnueli2001}
A.~Pnueli, S.~Ruah, and L.~Zuck, ``Automatic Deductive Verification
with Invisible Invariants,'' TACAS 2001.

\end{thebibliography}

\end{document}